%%%%%%%%%%%%%%%%%%%%%%%%%%%%%%%%%%%%%%%%%%%%%%%%%%%%%%%%%%%%%%%%%%%%%%%%%%%%%%%%
%%%%%%%%%%%%%%%%%%%%%%%%%%%%%%%%%%%%%%%%%%%%%%%%%%%%%%%%%%%%%%%%%%%%%%%%%%%%%%%%
%%                                                                            %%
%% thesistemplate.tex version 4.01 (2023/09/21)                               %%
%% The LaTeX template file to be used with the aaltothesis.sty (version 4.00) %%
%% style file.                                                                %%
%% This package requires pdfx.sty v. 1.5.84 (2017/05/18) or newer.            %%
%%                                                                            %%
%% This is licensed under the terms of the MIT license below.                 %%
%%                                                                            %%
%% Written by Luis R.J. Costa.                                                %%
%% Currently developed at Teacher services, Learning Services of Aalto        %%
%% University by Luis R.J. Costa since May 2019.                              %%
%%                                                                            %%
%% Copyright 2017-2021 aaltothesis.cls by Luis R.J. Costa,                    %%
%% luis.costa@aalto.fi.                                                       %%
%% Copyright 2017-2018 Swedish translations in aaltothesis.cls by Elisabeth   %%
%% Nyberg and Henrik Wallén henrik.wallen@aalto.fi.                           %%
%% Finnish documentation in the template opinnatepohja.tex translated from    %%
%% the English template documentation.                                        %%
%% Copyright 2021 English template thesistemplate.tex by Luis R.J. Costa,     %%
%% Maurice Forget, Henrik Wallén.                                             %%
%% Copyright 2018-2022 Swedish template kandidatarbetsbotten.tex by Henrik    %%
%% Wallen.                                                                    %%
%%                                                                            %%
%% Permission is hereby granted, free of charge, to any person obtaining a    %%
%% copy of this software and associated documentation files (the "Software"), %%
%% to deal in the Software without restriction, including without limitation  %%
%% the rights to use, copy, modify, merge, publish, distribute, sublicense,   %%
%% and/or sell copies of the Software, and to permit persons to whom the      %%
%% Software is furnished to do so, subject to the following conditions:       %%
%% The above copyright notice and this permission notice shall be included in %%
%% all copies or substantial portions of the Software.                        %%
%% THE SOFTWARE IS PROVIDED "AS IS", WITHOUT WARRANTY OF ANY KIND, EXPRESS OR %%
%% IMPLIED, INCLUDING BUT NOT LIMITED TO THE WARRANTIES OF MERCHANTABILITY,   %%
%% FITNESS FOR A PARTICULAR PURPOSE AND NONINFRINGEMENT. IN NO EVENT SHALL    %%
%% THE AUTHORS OR COPYRIGHT HOLDERS BE LIABLE FOR ANY CLAIM, DAMAGES OR OTHER %%
%% LIABILITY, WHETHER IN AN ACTION OF CONTRACT, TORT OR OTHERWISE, ARISING    %%
%% FROM, OUT OF OR IN CONNECTION WITH THE SOFTWARE OR THE USE OR OTHER        %%
%% DEALINGS IN THE SOFTWARE.                                                  %%
%%                                                                            %%
%%                                                                            %%
%%%%%%%%%%%%%%%%%%%%%%%%%%%%%%%%%%%%%%%%%%%%%%%%%%%%%%%%%%%%%%%%%%%%%%%%%%%%%%%%
%%                                                                            %%
%%                                                                            %%
%% An example for writting your thesis using LaTeX                            %%
%% Original version and development work by Luis Costa, changes to the text   %% 
%% in the Finnish template by Perttu Puska.                                   %%
%% Support for Swedish added 15092014                                         %%
%% PDF/A-b support added on 15092017                                          %%
%% PDF/A-2 support added on 24042018                                          %%
%% Layout design and typesettin changed 15072021                              %%
%%                                                                            %%
%% This example consists of the files                                         %%
%%       thesistemplate.tex (version 4.00) (for text in English)              %%
%%       opinnaytepohja.tex (version 4.00) (for text in Finnish)              %%
%%       kandidatarbetsbotten.tex (version 1.10) (for text in Swedish)        %%
%%       aaltothesis.cls                                                      %%
%%       linediagram.pdf (graphics file)                                      %%
%%       curves.pdf      (graphics file)                                      %%
%%       ledspole.jpg    (graphics file)                                      %%
%%                                                                            %%
%%                                                                            %%
%% Typeset in Linux with                                                      %%
%% pdflatex: (recommended method)                                             %%
%%             $ pdflatex thesistemplate                                      %%
%%             $ pdflatex thesistemplate                                      %%
%%                                                                            %%
%%   The result is the file thesistemplate.pdf that is PDF/A compliant, if    %%
%%   you have chosen the proper \documenclass options (see comments below)    %%
%%   and your included graphics files have no problems.                       %%
%%                                                                            %%
%%                                                                            %%
%% Explanatory comments in this example begin with the characters %%, and     %%
%% changes that the user can make with the character %                        %%
%%                                                                            %%
%%%%%%%%%%%%%%%%%%%%%%%%%%%%%%%%%%%%%%%%%%%%%%%%%%%%%%%%%%%%%%%%%%%%%%%%%%%%%%%%
%%%%%%%%%%%%%%%%%%%%%%%%%%%%%%%%%%%%%%%%%%%%%%%%%%%%%%%%%%%%%%%%%%%%%%%%%%%%%%%%
%%
%% WHAT is PDF/A
%%
%% PDF/A is the ISO-standardized version of the pdf. The standard's goal is to
%% ensure that he file is reproducable even after a long time. PDF/A differs
%% from pdf in that it allows only those pdf features that support long-term
%% archiving of a file. For example, PDF/A requires that all used fonts are
%% embedded in the file, whereas a normal pdf can contain only a link to the
%% fonts in the system of the reader of the file. PDF/A also requires, among
%% other things, data on colour definition and the encryption used.
%% Currently three PDF/A standards exist:
%% PDF/A-1: based on PDF 1.4, standard ISO19005-1, published in 2005.
%%          Includes all the requirements essential for long-term archiving.
%% PDF/A-2: based on PDF 1.7, standard ISO19005-2, published in 2011.
%%          In addition to the above, it supports embedding of OpenType fonts,
%%          transparency in the colour definition and digital signatures.
%% PDF/A-3: based on PDF 1.7, standard ISO19005-3, published in 2012.
%%          Differs from the above only in that it allows embedding of files in
%%          any format (e.g., xml, csv, cad, spreadsheet or wordprocessing
%%          formats) into the pdf file.
%% PDF/A-4: based on PDF 2.0, standard ISO19005-4, published in November 2020.
%%
%% PDF/A-1 files are not necessarily PDF/A-2 -compatible and PDF/A-2 are not
%% necessarily PDF/A-1 -compatible.
%% Standards PDF/A-1, PDF/A-2 and PDF/A-3 have two levels:
%% b: (basic) requires that the visual appearance of the document is reliably
%%    reproduceable.
%% a (accessible) in addition to the b-level requirements, specifies how
%%   accessible the pdf file is to assistive software, say, for the physically
%%   impaired.
%% The PDF/A-4 standard does not have additional levels like in the earlier
%% standards.
%% For more details on PDF/A, see, e.g., 
%% https://www.loc.gov/preservation/digital/formats/fdd/fdd000318.shtml or
%% https://www.pdfa.org/resource/iso-19005-pdfa/
%%
%%
%% WHICH PDF/A standard should my thesis conform to?
%%
%% Either to the PDF/A-1b or the PDF/A-2b standard. If all the figures and
%% graphs used in thesis work do not require transparency features, use either
%% PDF/A-1b or PFDF/A-2b. If you have figures with transparency
%% characteristics, use the PDF/A-2b standard. However, drawing applications
%% often use the transparency parameter, setting it to zero, to specify opacity
%% and get the basic 2-D visualisation. As a result, validation of PDF/A-1b
%% will fail. Use PDF/A-2b if PDF/A-1b validation fails.
%% Do not use the PDF/A-3b standard for your thesis.
%% The font to be used are specified in this templatenand they should not be
%% changed. In addition to not adhering to Aalto's visual guidelines, you may
%% have difficulties in producing a PDF/A-compliant PDF.
%%
%%
%% Validate your PDF/A file at https://www.pdf-online.com/osa/validate.aspx
%%
%%
%% WHAT graphics format can I use to produce my PDF/A compliant file?
%%
%% When using pdflatex to compile your work, favour the use of pdf, but you can
%% use the jpg or png format especially for photographs. You will have PDF/A 
%% compliance problems with figures in pdf if the fonts are not embedded in the
%% pdf file.
%% If you choose to use latex to compile your work, the only acceptable file
%% format for your figure is eps. DO NOT use the ps format for your figures.

%% USE one of the following three \documentclass set-ups:
%% * the first when using pdflatex to directly typeset your document in the
%%   chosen pdf/a format for online publishing (centred page layout),
%% * the second for one-sided printing your thesis with the layout (wide left 
%%   margin), or
%% * the third for two-sided printing.
%%
\PassOptionsToPackage{main=english,finnish,swedish}{babel}
\documentclass[english, 12pt, a4paper, sci, utf8, a-2b, online]{aaltothesis}
%\documentclass[english, 12pt, a4paper, elec, utf8, a-2b, print]{aaltothesis}
%\documentclass[english, 12pt, a4paper, elec, utf8, a-2b, print, twoside]{aaltothesis}

%% Use the following options in the \documentclass macro above:
%% your school: arts, biz, chem, elec, eng, sci
%% the character encoding scheme used by your editor: utf8, latin1
%% thesis language: english, finnish, swedish
%% make an archiveable PDF/A-1b or PDF/A-2b compliant file: a-1b, a-2b
%%                    (with pdflatex, a normal pdf containing metadata is
%%                     produced without the a-*b option)
%% typset for online document or print on paper: online, print
%%        online: typeset in symmetric layout and blue hypertext for online
%%                publishing
%%        print: typeset in a symmetric layout and black hypertext for printing
%%               on paper
%%          two-side printing: twoside (default is one-sided printing)
%%               typeset in a wide margin on the binding side of the page and
%%               black hypertext. Use with print only.
%%

%% Use one of these if you write in Finnish (or use the Finnish template
%% opinnaytepohja.tex)
%\documentclass[finnish, 12pt, a4paper, elec, utf8, a-1b, online]{aaltothesis}
%\documentclass[finnish, 12pt, a4paper, elec, utf8, a-1b, print]{aaltothesis}
%\documentclass[finnish, 12pt, a4paper, elec, utf8, a-1b, print, twoside]{aaltothesis}

%% Use one of these if you write in Swedish (or use the Swedish template
%% kandidatarbetsbotten.tex)
%\documentclass[swedish, 12pt, a4paper, elec, utf8, a-2b, online]{aaltothesis}
%\documentclass[swedish, 12pt, a4paper, elec, utf8, a-2b]{aaltothesis}
%\documentclass[swedish, 12pt, a4paper, elec, dvips, online]{aaltothesis}

%% FOR USERS OF AMS PACKAGES:
%% * newtxmath used in this template loads amsmath, so
%%   you needn't load it. If you want to use options in amsmath, load it here, 
%%   before \setupthesisfonts below to pass the options to amsmath.
%% * If you want to use amsthm, load it here before \setupthesisfonts to avoid
%%   a clash with newtxmath.
%% * If using amsmath with options and you want to use amsthm, load amsthms
%%   after amsmath, as described in the amsthm documentation.
%% * Don't use amsbsym or amsfonts. The symbols [and macros] there are defined in
%%   newtxmath and so clash if used.
%\usepackage[options]{amsmath}
%\usepackage{amsthm}

%% DO NOT MOVE OR REMOVE \setupthesisfonts
\setupthesisfonts

%%
%% Add here the packges you need
%%
\usepackage{graphicx}
\graphicspath{{figures/}{../docs/figures/}{../docs/architectures/}}
\usepackage{subfig}
\usepackage{natbib}
\usepackage{dsfont}
\usepackage{amsmath}
\usepackage{mathtools}
\usepackage{algpseudocode}
\usepackage{algcompatible}
\usepackage{makecell}
%% For tables that span multiple pages; used to split a paraphrasing example in
%% the appendix. If you don't need it, remove it.
\usepackage{longtable}
\usepackage{booktabs}
\usepackage{placeins}

%% A package for generating Creative Commons copyright terms. If you don't use
%% the CC copyright terms, remove it, since otherwise undesired information may
%% be added to this document's metadata.
\usepackage[type={CC}, modifier={by-nc-sa}, version={4.0}]{doclicense}
%% Find below three examples for typesetting the CC license notice.


%% Edit to conform to your degree programme
%% Capitalise the words in the name of the degree programme: it's a name
\degreeprogram{Computer, Communication and Information Sciences}
%%

%% Your major
%%
\major{Machine Learning, Data Science and Artificial Intelligence}
%%

%% Choose one of the three below
%%
%\univdegree{BSc}
\univdegree{MSc}
%\univdegree{Lic}
%%

%% Your name (self explanatory...)
%%
\thesisauthor{Reza Soumi}
%%

%% Your thesis title and possible subtitle comes here and possibly, again,
%% together with the Finnish or Swedish abstract. Do not hyphenate the title
%% (and subtitle), and avoid writing too long a title. Should LaTeX typeset a
%% long title (and/or subtitle) unsatisfactorily, you might have to force a
%% linebreak using the \\ control characters. In this case...
%% * Remember, the title should not be hyphenated!
%% * A possible 'and' in the title should not be the last word in the line; it
%%   begins the next line.
%% * Specify the title (and/or subtitle) again without the linebreak characters
%%   in the optional argument in box brackets. This is done because the title
%%   is part of the metadata in the pdf/a file, and the metadata cannot contain
%%   linebreaks.
%%
\thesistitle[Learning Shared and Private Representations for Multi-omics Cancer Subtyping]{Learning Shared and Private\\ Representations for Multi-omics\\ Cancer Subtyping}
%\thesistitle[Title of the thesis]{Title of\\ the thesis}
%%
%% Either remove or leave \thesissubtitle{} empty if you don't use it
%%
% \thesissubtitle{A possible subtitle}
%\thesissubtitle[Subtitle of the thesis]{Subtitle of\\ the thesis}
%\thesissubtitle{}

%%
\place{Espoo}
%%

%% The date for the bachelor's thesis is the day it is presented
%%
\date{8 November 2026}
%%

%% Thesis supervisor
%% Note the "\" character in the title after the period and before the space
%% and the following character string.
%% This is because the period is not the end of a sentence after which a
%% slightly longer space follows, but what is desired is a regular interword
%% space.
%%
\supervisor{Prof.\ Pekka Marttinen}
%%

%% Advisor(s)---two at the most---of the thesis. Check with your supervisor how
%% many official advisors you can have.
%%
\advisor{Negar Safinianaini (PhD)}
%%

%% If you do your thesis work in a company of other institute, give the name of
%% the company or instution here. Otherwise, leave the macro empty, comment it
%% out, or remove it. This will remove this field from the abstract page.
%%
% \collaborativepartner{Company or institute name (if relevant)}
%%

%% Aaltologo: syntax:
%% \uselogo{?|!|''}
%% The logo language is set to be the same as the thesis language.
%%
%\uselogo{?}
%\uselogo{!}
\uselogo{''}
%%

%%%%%%%%%%%%%%%%%%               COPYRIGHT TEXT               %%%%%%%%%%%%%%%%%%
%%%%%%%%%%%%%%%%%%%%%%%%%%%%%%%%%%%%%%%%%%%%%%%%%%%%%%%%%%%%%%%%%%%%%%%%%%%%%%%%

%% Copyright of a work is with the creator/author of the work regardless of
%% whether the copyright mark is explicitly in the work or not. You may, if you
%% wish---we encourage you to do so---publish your work under a Creative
%% Commons license (see creativecommons.org), in which case the license text
%% must be visible in the work. Write here the copyright text you want using the
%% macro \copyrighttext, which writes the text into the metadata of the pdf file
%% as well.
%%
%% Syntax:
%% \copyrigthtext{metadata text}{text visible on the page}
%%
%% CHOOSE ONE OF THE COPYRIGHT NOTICE STYLES BELOW.
%% IF USING THE CC TERMS, CHOOSE THE LICENSE YOU WANT TO USE.
%% The different CC licenses are listed at 
%% https://creativecommons.org/about/cclicenses/.
%% If you use the icons from the dolicense.sty package, add the package above
%% (\usepackage{dolicense}).
%% IMPORTANT NOTE!! Manually write the CC text in the \copyrighttext metadata
%% text field.
%%
%% NOTE: In the macros below, the text written in the metadata must have a
%% \noexpand macro before the \copyright special character. When not in pdf/a
%% mode (i.e. a-1b or a-2b are not specified in \documentclass), two \noexpands
%% are required in the metadata text to correctly render the copyright mark in
%% the pdf metadata. In pdf/a mode one \noexpand suffices.
%%
%% EXAMPLE OF PLAIN COPYRIGHT TEXT
%% The macros \copyright and \year below must be separated by the \ character 
%% (space chacter) from the text that follows. The macros in the argument of the
%% \copyrighttext macro automatically insert the year and the author's name.
%% (Note! \ThesisAuthor is an internal macro of the aaltothesis.cls class file).
%%
%\copyrighttext{Copyright \noexpand\textcopyright\ \number\year\ \ThesisAuthor}
%{Copyright \textcopyright{} \number\year{} \ThesisAuthor}
%%
%% Of course, the same text could have simply been written as
%% \copyrighttext{Copyright \noexpand\copyright\ 2018 Eddie Engineer}
%% {Copyright \copyright{} 2022 Eddie Engineer}
%%
%% EXAMPLES OF CC LICENSE: different ways to display the same license
%% 1. A simple Creative Commons license text with a link to the copyright notice:
%\copyrighttext{\noexpand\textcopyright\ \number\year. This work is 
%	licensed under a CC BY-NC-SA 4.0 license.}{\textcopyright{} 
%	\number\year. This work is licensed under a 
%	\href{https://creativecommons.org/licenses/by-nc-nd/4.0/}{CC BY-NC-SA 4.0} 
%	license.}
%
%% To get the URL of the license of your choice, go to 
%% https://creativecommons.org/about/cclicenses/, click on the chosen license
%% you want to use, and copy-and-paste the URL in the macro \href above.
%%
%% 2. A short Creative Commons license text containing the respective CC icons
%% (requires the package dolicense.sty to be added in the preamble as done
%% above) and a link to the corresponding Creative Commons license webpage (see
%% the dolicense package documentation for other license icons):
%\copyrighttext{\noexpand\textcopyright\ \number\year. This work is licensed
%	under a CC BY-NC-SA 4.0 license.}{
%	\parbox{95mm}{\noindent\textcopyright\ \number\year. \doclicenseText} 
%	\hspace{1em}\parbox{35mm}{\doclicenseImage}
%}
%%
%% 3. An expanded Creative Commons license text containing the respective CC
%% icons text and as generated by the dolicense.sty package (the license is set
%% via package options in \usepackage[options]{dolicense} above; see the
%% dolicense package documentation for other license texts and icons):
\copyrighttext{\noexpand\textcopyright\ \number\year. This work is 
	licensed under a Creative Commons "Attribution-NonCommercial-ShareAlike 4.0 
	International" (BY-NC-SA 4.0) license.}{\noindent\textcopyright\ \number
	\year \ \doclicenseThis}
%%%%%%%%%%%%%%%%%%%%%%%%%%%%%%%%%%%%%%%%%%%%%%%%%%%%%%%%%%%%%%%%%%%%%%%%%%%%%%%%


%% The English abstract:
%% All the details (name, title, etc.) on the abstract page appear as specified
%% above.
%% Thesis keywords:
%% Note! The keywords are separated using the \spc macro
%%
\keywords{multi-omics\spc cancer subtyping\spc shared representation\spc private representation\spc disentanglement}
%%

%% The abstract text. This text in one paragraph is included in the metadata of
%% the pdf file as well as the abstract page. To have paragraphs in your
%% abstract rewrite it in the abstarct environment as described below.
%%
\thesisabstract{%
Deep multi-omics models often learn shared and private latent codes for cancer
subtyping and report strong unsupervised clustering metrics. This thesis asks
whether those codes actually play their intended roles. On TCGA breast cancer
data with multi-seed evaluation, MOCSS yields high integration NMI while private
codes retain subtype structure and cross-view reconstruction gaps stay near zero.
OMIDIENT appears better separated on private-subtype NMI, but probes indicate
private-path collapse rather than verified privacy. Three methods are developed
to enforce a clearer split: cross-view predictability constraints, joint
reconstruction with VICReg independence, and spectral alignment of shared codes
to a fused multi-omics patient graph. Results show that low private-subtype NMI
is not a reliable disentanglement trophy, that Method 2 is preferable to the
VICReg ablation of Method 1 for integration quality, and that Method 3 achieves
the strongest cross-view shared-versus-private distinction among AE-private
models while remaining competitive on clustering.
}

%% You can prevent LaTeX from writing into the xmpdata file (it contains all the 
%% metadata to be written into the pdf file) by setting the writexmpdata switch
%% to 'false'. This allows you to write the metadata in the correct format
%% directly into the file thesistemplate.xmpdata.
%\setboolean{writexmpdatafile}{false}


%% All that is printed on paper starts here
%%
\begin{document}

%% Create the coverpage
%%
\makecoverpage

%% Typeset the copyright text.
%% If you wish, you may leave out the copyright text from the human-readable
%% page of the pdf file. This may seem like a attractive idea for the printed
%% document especially if "Copyright (c) yyyy Eddie Engineer" is the only text
%% on the page. However, the recommendation is to print this copyright text.
%%
\makecopyrightpage

\clearpage
%% Note that when writing your thesis in English, place the English abstract
%% first followed by the possible Finnish or Swedish abstract.

%% Abstract text
%% All the details (name, title, etc.) on the abstract page appear as specified
%% above. Add your abstarct text with paragraphs here to have paragraphs in the
%% visible abstract page. Nonetheless, write the abstarct text without
%% paragraphs in the macro \thesismacro so that it is added to the metadata as
%% well.
%%
\begin{abstractpage}[english]
  Deep multi-omics models often learn shared and private latent codes for cancer
  subtyping and report strong unsupervised clustering metrics. This thesis asks
  whether those codes actually play their intended roles. On TCGA breast cancer
  data with multi-seed evaluation, MOCSS yields high integration NMI while private
  codes retain subtype structure and cross-view reconstruction gaps stay near zero.
  OMIDIENT appears better separated on private-subtype NMI, but probes indicate
  private-path collapse rather than verified privacy. Three methods are developed
  to enforce a clearer split: cross-view predictability constraints, joint
  reconstruction with VICReg independence, and spectral alignment of shared codes
  to a fused multi-omics patient graph. Results show that low private-subtype NMI
  is not a reliable disentanglement trophy, that Method~2 is preferable to the
  VICReg ablation of Method~1 for integration quality, and that Method~3 achieves
  the strongest cross-view shared-versus-private distinction among AE-private
  models while remaining competitive on clustering.

  % The abstract is a short description of the essential contents of the thesis,
  % usually in one paragraph: what was studied and how and what were the main
  % findings.

  % For a Finnish thesis, the abstract should be written in both Finnish and
  % English; for a Swedish thesis, in Swedish and English. The abstracts for
  % English theses written by Finnish or Swedish speakers should be written in
  % English and either in Finnish or in Swedish, depending on the student’s
  % language of basic education. Students educated in languages other than Finnish
  % or Swedish write the abstract only in English. Students may include a second
  % or third abstract in their native language, if they wish.

  % The abstract text of this thesis is written on the readable abstract page as
  % well as into the pdf file's metadata via the \verb+\thesisabstract+ macro
  % (see comment in this \TeX{} file above). Write here the text that goes onto
  % the readable abstract page. You can have special characters, linebreaks, and
  % paragraphs here. Otherwise, this abstract text must be identical to the
  % metadata abstract text.
  
  % If your abstract does not contain special characters and it does not require
  % paragraphs, you may take advantage of the \verb+\abstracttext+ macro (see the
  % comment in this \TeX{} file below).
\end{abstractpage}

%% The text in the \thesisabstract macro is stored in the macro \abstractext, so
%% you can use the text metadata abstract directly as follows:
%%
%\begin{abstractpage}[english]
%	\abstracttext{}
%\end{abstractpage}

%% Force a new page so that the possible Finnish or Swedish abstract does not
%% begin on the same page
%%
\newpage

\dothesispagenumbering{}

%% Front matter: preface, table of contents, symbols and abbreviations
%% Included from thesistemplate.tex (no \documentclass / preamble here).

\mysection{Preface}

This Master's thesis was carried out in the Machine Learning for Health research group at the Department of Computer Science, Aalto University. I am grateful to Professor Pekka Marttinen for supervising the work and for continuous scientific guidance on representation learning and probabilistic modelling. I thank Negar Safinianaini for co-supervision, for introducing me to the OMIDIENT line of research, and for many detailed discussions on shared--private multi-omics models and evaluation practice. Their feedback substantially shaped both the methodological choices and the critical interpretation of the experimental results.

I also thank colleagues and fellow students who provided informal feedback on figures, writing, and experimental design. Any remaining errors are my own.

\vspace{5cm}
Otaniemi, November 2026\\

\vspace{5mm}
{\hfill Reza Soumi \hspace{1cm}}

\newpage

\thesistableofcontents

\mysection{Symbols and abbreviations}

\subsection*{Symbols}

\begin{tabular}{ll}
$M$ & number of omics modalities (views) \\
$m$ & modality index, $m \in \{1,\ldots,M\}$ \\
$N$ & number of samples (patients) \\
$X_m$ & observed feature matrix of modality $m$ \\
$S_m$ & intended shared latent code of modality $m$ \\
$P_m$ & intended private / view-specific latent code of modality $m$ \\
$\bar{S}$ & mean shared embedding, $\bar{S}=\mathrm{mean}_m(S_m)$ \\
$[P;S]$ & concatenation of private and shared codes \\
$U$ & offline fused spectral target (Laplacian embedding) \\
$r$ & spectral target dimension (here $r=16$) \\
$\mathrm{Dec}$, $D$ & decoder / reconstruction network \\
$h,f,g,k,b$ & Method~1 predictability networks \\
$f_{\mathrm{spec}}$ & spectral regression head mapping $\bar{S}$ to $U$ \\
$\lambda_{\mathrm{ctr}},\lambda_{\mathrm{vic}},\ldots$ & loss weights for contrastive, VICReg, and related terms \\
$\Delta\mathrm{X}\,R^{2}$ & cross-view probe gap $R^{2}(X\leftarrow S_{\mathrm{other}})-R^{2}(X\leftarrow P_{\mathrm{other}})$ \\
$\mathrm{NMI}$ & normalised mutual information \\
$\mathrm{ARI}$ & adjusted Rand index \\
$\mathrm{CKA}$ & centred kernel alignment \\
$\mathcal{L}$ & training objective (loss) \\
\end{tabular}

\subsection*{Operators}

\begin{tabular}{ll}
$\|\,\cdot\,\|$ & matrix or vector norm (as indicated in context) \\
$\mathrm{concat}(\cdot)$ & concatenation of latent codes \\
$\mathrm{stopgrad}(\cdot)$ & stop-gradient (detach) operator \\
$\mathbb{E}[\,\cdot\,]$ & expectation \\
$R^{2}(Y\leftarrow Z)$ & coefficient of determination of Ridge regression from $Z$ to $Y$ \\
\end{tabular}

\subsection*{Abbreviations}

\begin{tabular}{ll}
AE & autoencoder \\
ARI & adjusted Rand index \\
BRCA & breast invasive carcinoma (TCGA cohort) \\
CCA & canonical correlation analysis \\
CKA & centred kernel alignment \\
GRL & gradient reversal layer \\
InfoNCE & information noise-contrastive estimation (contrastive loss) \\
JIVE & Joint and Individual Variation Explained \\
KIRC & kidney renal clear cell carcinoma (TCGA cohort) \\
MLP & multilayer perceptron \\
MOCSS & Multi-Omics Contrastive Shared-Specific learning \\
MOFA & Multi-Omics Factor Analysis \\
NMI & normalised mutual information \\
OMIDIENT & multi-omics integration with Dirichlet autoencoder networks \\
PAM50 & Prediction Analysis of Microarray 50 (breast cancer subtypes) \\
SNF & Similarity Network Fusion \\
TCGA & The Cancer Genome Atlas \\
UMAP & Uniform Manifold Approximation and Projection \\
VAE & variational autoencoder \\
VICReg & Variance--Invariance--Covariance Regularization \\
\end{tabular}


\cleardoublepage

\section{Introduction}

Cancer is a molecularly heterogeneous disease. Patients who share a clinical diagnosis often differ in transcriptional programmes, epigenetic regulation, and post-transcriptional control, and those differences matter for prognosis and therapy selection~\citep{hasin2017multi,jiang2022big}. Multi-omics profiling---joint measurement of several molecular layers for the same samples---has therefore become a central resource for unsupervised discovery of tumour subtypes~\citep{mohr2024navigating}. Among public resources, The Cancer Genome Atlas (TCGA) remains the most widely used benchmark for method development, with breast invasive carcinoma (BRCA) providing a particularly mature subtype ontology through the PAM50 labels~\citep{parker2009pam50}.

Integrating heterogeneous omics views is not merely a matter of concatenating features. Modalities differ in dimension, noise structure, and biological content: some signals are shared across assays, while others are view-specific. Classical integrative methods such as iCluster~\citep{shen2009integrative}, Multi-Omics Factor Analysis (MOFA)~\citep{argelaguet2018multi}, Joint and Individual Variation Explained (JIVE)~\citep{lock2013jive}, and Similarity Network Fusion (SNF)~\citep{wang2014snf} already encode this intuition by separating joint from individual structure, or by fusing patient similarity graphs. Deep learning has extended the same idea to nonlinear latent spaces. Shared--private autoencoders learn, for each modality~$m$, a shared code~$S_m$ intended to capture cross-view biology and a private code~$P_m$ intended to capture residual view-specific variation~\citep{chen2023mocss,xiao2022multi}. Contrastive alignment of projected shared codes, as in MOCSS~\citep{chen2023mocss}, encourages cross-omics consistency, while OMIDIENT replaces the private branch with a product-Dirichlet variational model while retaining a MOCSS-style shared pathway~\citep{Safinianaini_Välimäki_Bresson_Gorbonos_Rajamäki_Aaltonen_Marttinen_2025}.

The practical success of these models is usually summarised by unsupervised clustering metrics on an integrated embedding---most often normalised mutual information (NMI) with clinical subtypes. High NMI is valuable, but it does not certify that~$S$ and~$P$ play their intended roles. A private code that quietly absorbs subtype-discriminative shared biology, or a shared code that mainly reconstructs local assay detail, can still yield competitive clustering after concatenation. Orthogonality penalties that act only through weak elementwise products do not guarantee statistical independence, and dual reconstruction paths that require both~$S_m$ and~$P_m$ to reconstruct the full input~$X_m$ actively encourage redundancy. Representation-learning research outside multi-omics has repeatedly shown that reconstruction alone does not identify the factors one cares about~\citep{bengio2013representation,Goodfellow-et-al-2016}; the same warning applies here.

\paragraph{Motivating observation.}
Preliminary diagnostics on BRCA, including embedding visualisations of the kind later summarised in Figure~\ref{fig:umap} (Chapter~\ref{sec:results}), suggest that models with similar integration NMI can exhibit qualitatively different shared--private geometries. In MOCSS, private embeddings often retain clear subtype structure, while the shared codes look only partially aligned with the same taxonomy. In OMIDIENT, private codes appear nearly non-informative for subtypes under concatenation metrics, yet own-view reconstruction probes indicate pathological collapse rather than clean privacy. Quantitative companions of those figures---NMI bars, metric heatmaps, and scatter plots of subtype separation against~$\Delta\mathrm{X}\,R^{2}$---point in the same direction. These observations motivate an evaluation protocol that treats role separation as a first-class scientific claim, not as a corollary of clustering accuracy.

The thesis therefore separates two questions that are often conflated. The \emph{clustering} question asks whether an embedding recovers PAM50 structure. The \emph{decomposition} question asks whether cross-view signal concentrates in~$S$ while~$P$ remains residual or complementary. Answering the second requires probes that clustering tables omit.

\subsection{Research questions}

This thesis addresses the following questions.

\begin{enumerate}
\item Do high unsupervised subtype NMI scores of prior shared--private models (MOCSS, OMIDIENT) imply that~$S$ and~$P$ carry the intended complementary information on BRCA multi-omics data?
\item Which label-free probes---in particular cross-view reconstruction gaps~$\Delta\mathrm{X}\,R^{2}$, own-view Ridge~$R^{2}$, centred kernel alignment (CKA), and embedding visualisations---best diagnose failures of shared--private decomposition?
\item Can architectural and objective changes that enforce a clearer distinction between~$S$ and~$P$ (Methods~1--3 introduced in this thesis) improve role separation without sacrificing competitive integration performance?
\item Among the proposed methods, which trade-offs between predictability constraints, VICReg-style independence, joint reconstruction, and spectral alignment are empirically preferable under a multi-seed BRCA evaluation?
\end{enumerate}

\subsection{Objectives and scope}

The objective is not to maximise subtype NMI at all costs. Instead, the thesis aims to (i)~audit existing shared--private pipelines with probes that test whether foreign private codes predict another view better or worse than foreign shared codes; (ii)~design Methods~1--3 that operationalise shared content as cross-view predictable structure and private content as residual, view-local variation; and (iii)~compare methods under a fixed encoder backbone, label-free model selection, and multi-seed reporting. The empirical focus is TCGA BRCA with mRNA, DNA methylation, and miRNA views and PAM50 subtypes. Kidney renal clear cell carcinoma (KIRC) is included in the experimental setup as a secondary cohort; KIRC numerical results will be reported in a later revision of this thesis. Method~4 from the accompanying codebase is out of scope and is not discussed.

\subsection{Contributions}

The main contributions are as follows.

\begin{enumerate}
\item A critical multi-seed audit of MOCSS and OMIDIENT on BRCA showing that competitive integration NMI coexists with failed shared--private role assignment: MOCSS private codes steal subtype signal, while OMIDIENT private codes collapse under own-view probes despite large~$\Delta\mathrm{X}\,R^{2}$.
\item An evaluation protocol centred on~$\Delta\mathrm{X}\,R^{2}$, role-aware UMAP inspection, and integration NMI, with an explicit rule against ranking models by minimising $\mathrm{concat}(P)$ subtype NMI alone.
\item Three methods designed to enforce shared--private distinction: Method~1 (cross-view predictability with adversaries and private sufficiency), a Method~1+VICReg ablation, Method~2 (joint reconstruction with VICReg and contrastive alignment), and Method~3 (offline fused spectral targets with residual private decoding).
\item A BRCA multi-seed ($n=5$, seeds $\{21,\ldots,25\}$) empirical comparison identifying Method~3 as the clearest winner on~$\Delta\mathrm{X}\,R^{2}$ while remaining competitive on concatenation NMI, and Method~2 as a simpler AE-private alternative preferred over the VICReg ablation when both role separation and integration are considered jointly.
\end{enumerate}

\subsection{Structure of the thesis}

Chapter~\ref{sec:literature} reviews classical and deep multi-omics integration, with emphasis on shared--private models and evaluation pitfalls. Chapter~\ref{sec:background} summarises the biological setting and the technical background on autoencoders, contrastive learning, VICReg, spectral clustering, and clustering metrics. Chapter~\ref{sec:methodology} formalises the problem, describes baselines and Methods~1--3, and defines the evaluation protocol. Chapter~\ref{sec:results} reports BRCA multi-seed results, leading with shared--private diagnostics before subtyping performance. Chapter~\ref{sec:discussion} interprets the findings, limitations, and implications for future multi-omics representation learning. Chapter~\ref{sec:conclusions} concludes.

\clearpage
\section{Literature review}
\label{sec:literature}

Multi-omics integration sits at the intersection of computational biology and representation learning. This chapter reviews classical statistical approaches, deep multi-omics models, the shared--private contrastive line of work that includes MOCSS and OMIDIENT, and broader machine-learning critiques of claimed disentanglement. The chapter closes by stating the gap addressed in this thesis.

\subsection{Classical integrative clustering and factorisation}

Early multi-omics subtype discovery largely relied on linear latent-variable or graph-fusion models. iCluster~\citep{shen2009integrative} poses integrative clustering as a joint Gaussian latent-variable model: a common discrete cluster assignment generates modality-specific continuous loadings, so that subtype structure is forced to be shared while residual variation remains view-local. The method is statistically clear and remains a common baseline, but scalability and nonlinear interactions are limited.

MOFA~\citep{argelaguet2018multi} and related factor models extend the idea to continuous latent factors with sparsity and likelihood choices adapted to each omics type. Factors can be shared or private depending on which views they load on, providing an interpretable decomposition without requiring a single clustering objective. JIVE~\citep{lock2013jive} makes the shared--individual partition explicit at the matrix level by decomposing a multi-block data set into joint structure, individual structure, and residual noise. These methods establish the conceptual vocabulary used throughout this thesis: shared structure should explain cross-view agreement, while private structure should explain leftover within-view variance.

Graph-based approaches take a complementary route. Similarity Network Fusion~\citep{wang2014snf} constructs a patient similarity network per modality and iteratively fuses them so that local neighbourhoods that agree across views are reinforced. Spectral clustering on the fused graph then yields subtypes~\citep{ng2002spectral,vonluxburg2007tutorial}. SNF does not learn parametric encoders, but it demonstrates that cross-view agreement is naturally expressed as geometry on samples rather than as independent per-view reconstructions. That geometric intuition reappears in Method~3 of this thesis.

Benchmarking studies have emphasised that no single classical method dominates all cancers and metrics~\citep{rappoport2021evaluation,cantini2021benchmarking}. Performance depends on preprocessing, the number of clusters, and whether the biological ground truth is itself contested. Consequently, modern work increasingly reports multiple metrics and stress-tests robustness across seeds and hyperparameter choices~\citep{baiao2025technical}.

A further lesson from the classical literature is that ``shared'' has always been an operational concept. In iCluster it is the discrete assignment; in MOFA and JIVE it is the set of factors or matrix components with multi-view support; in SNF it is the fused neighbourhood relation. None of these definitions reduces to ``whatever remains after training an autoencoder with an orthogonality penalty.'' That historical precision motivates the probe-centred evaluation used later in this thesis.

\subsection{Deep multi-omics representation learning}

Deep models replace linear factors with nonlinear encoders, typically autoencoders or variational autoencoders~\citep{kingma2013auto,Goodfellow-et-al-2016}. The appeal is capacity: nonlinear maps can capture interactions among genes, CpG sites, and microRNAs that linear factor models miss. Surveys of multi-omics machine learning highlight both the promise and the risk of opaque embeddings that are hard to validate biologically~\citep{hasin2017multi,xiao2022multi,jiang2022big,mohr2024navigating}. Large-scale cancer genomics programmes have accelerated the availability of matched multi-omics, but they have also increased the temptation to treat any embedding that clusters by clinical labels as scientifically resolved~\citep{jiang2022big}.

Architecturally, deep integration strategies fall into early, intermediate, and late fusion. Intermediate fusion---learning per-view encoders and then combining latent codes---is the dominant pattern for unsupervised subtyping because it respects modality heterogeneity while still enabling joint clustering. Early concatenation ignores scale and noise differences; late fusion of separately clustered views can miss coordinated but weak cross-omics signal. Intermediate shared--private models attempt to keep the best of both worlds by allocating capacity to cross-view and view-local factors inside a common training loop~\citep{xiao2022multi}.

Contrastive learning has become a popular alignment tool in this setting: positive pairs are matched samples across views, and InfoNCE-style objectives pull their embeddings together~\citep{chen2020simple}. When applied to multi-omics, contrastive losses encourage cross-view consistency of the aligned subspace, but they do not by themselves define what should remain private. If the shared encoder can still reconstruct local assay detail, contrastive alignment may merely synchronise redundant copies of mixed information.

Several deep shared--private designs have been proposed. The common template encodes each view into~$S_m$ and~$P_m$, applies a cross-view consistency loss on~$S$, and reconstructs~$X_m$ from some combination of~$S_m$ and~$P_m$. The details of reconstruction and independence regularisation determine whether the intended semantics hold. Weak orthogonality between~$S_m$ and~$P_m$ is often used as a cheap surrogate for disentanglement, yet elementwise product penalties do not imply subspace orthogonality or statistical independence~\citep{zbontar2021barlow,bardes2022vicreg}. Principal-component diagnostics of omics matrices already show that dominant variance directions need not coincide with biologically annotated structure~\citep{yeung2001principal}; likewise, dominant reconstructive directions in a private code need not be ``private'' in the scientific sense.

\subsection{MOCSS and OMIDIENT}

MOCSS~\citep{chen2023mocss} is a representative deep shared--specific model for cancer subtyping. Each modality is mapped by an autoencoder into shared and specific latents. Two design choices are especially consequential. First, both the shared decoder and the specific decoder reconstruct the full input~$X_m$, so each code is pressured to be individually sufficient for reconstruction. Second, shared codes are projected by a multilayer perceptron and aligned with a contrastive (InfoNCE) loss across modalities, while a weak elementwise orthogonality term discourages co-activation of~$S_m$ and~$P_m$. Empirically, MOCSS reports strong unsupervised subtype recovery on TCGA cohorts. The present thesis revisits that success through role-aware probes and finds that high concatenation NMI can coexist with private codes that retain---and sometimes dominate---subtype structure. Equally telling is the cross-view reconstruction pattern: predicting a modality~$X_A$ from the shared codes of the other modalities yields approximately the same Ridge~$R^{2}$ as predicting~$X_A$ from those modalities' private codes, so the cross-view gap~$\Delta\mathrm{X}\,R^{2}$ stays near zero. If~$S$ and~$P$ truly played distinct roles, foreign shared codes should systematically outperform foreign private codes at explaining another view; the near-parity of the two predictors is further evidence that the named shared and private latents are not what the architecture labels suggest.

From a modelling standpoint, dual full reconstruction is in tension with the residual interpretation of private codes. If~$P_m$ must reconstruct~$X_m$ alone, it has little incentive to leave cross-view signal exclusively to~$S_m$. Orthogonality that penalises only per-dimension co-activation after normalisation cannot close that loophole: nonlinear mixing and subspace overlap remain possible. The contrastive term improves cross-view agreement of projected~$S$, but agreement is compatible with both pure shared biology and synchronised leakage of private detail.

OMIDIENT~\citep{Safinianaini_Välimäki_Bresson_Gorbonos_Rajamäki_Aaltonen_Marttinen_2025} retains a MOCSS-style shared pathway but replaces the private branch with a product-Dirichlet variational encoder--decoder. Its primary design goal is to improve the \emph{private/specific} representation---exploiting compositional Dirichlet structure and modelling view-specific uncertainty more faithfully than a deterministic autoencoder---and the original work reports gains on clustering, classification, and missing-data reconstruction that demonstrate the value of that private-branch change. The paper's central claim is therefore about a better private model within a shared--specific template, not about establishing a new operational definition of sharedness or about proving that~$S$ and~$P$ are information-theoretically distinct. In principle, a well-behaved private Dirichlet code should absorb residual assay-specific variation while leaving cross-view subtype biology in~$S$. When the same BRCA checkpoints are re-examined with the role probes of this thesis, multi-seed evaluation reveals near-zero $\mathrm{concat}(P)$ NMI together with strongly negative own-view $X\leftarrow P$ Ridge~$R^{2}$. That pattern is consistent with private-path collapse under the probe suite used here, rather than with a verified shared--private theory win. The point is not that OMIDIENT lacks scientific value---on the contrary, it advances private modelling and remains a strong integration baseline---but that low private-subtype NMI alone should not be read as proof that the shared/private labels have been validated. Moreover, single-seed subtype scores that look favourable can shrink under multi-seed aggregation, underscoring the need for seed-robust reporting~\citep{baiao2025technical}. Because OMIDIENT's shared branch is still MOCSS-like, improvements confined to the private likelihood do not automatically repair the definition of sharedness; that complementary question is the focus of the present thesis.

\subsection{Shared--private learning beyond multi-omics}

Outside genomics, representation learning has long distinguished factors that should be invariant across views from factors that should remain style- or domain-specific~\citep{bengio2013representation,Goodfellow-et-al-2016}. Contrastive methods such as SimCLR~\citep{chen2020simple} maximise agreement under augmentation, while non-contrastive methods such as Barlow Twins~\citep{zbontar2021barlow} and VICReg~\citep{bardes2022vicreg} control redundancy through covariance and variance constraints. These ideas are transferable to multi-omics if one treats modalities as views: shared codes should agree across views, and private codes should explain residuals without duplicating the shared subspace.

Critiques of claimed disentanglement emphasise that independence penalties, reconstruction losses, and clustering metrics measure different properties~\citep{martens2024disentangling}. A model can look disentangled under one probe and entangled under another. In multi-omics, an especially misleading failure mode is to treat low private-subtype NMI as proof of successful separation: private codes may be collapsed, or they may encode complementary biology that is subtype-related yet not redundant with~$S$. Levine et al.~\citep{levine2013integrated} and related clinical multi-omics studies illustrate that assay-specific correlates of subtype are scientifically meaningful; discarding them by force can harm both interpretation and integration. Navigating multi-omics method choice therefore requires matching the inductive bias to the scientific question, not only to a leaderboard metric~\citep{mohr2024navigating}.

\subsection{Evaluation practice and the research gap}

Technical surveys caution that preprocessing, feature selection, and metric choice can dominate apparent method rankings~\citep{baiao2025technical,rappoport2021evaluation,cantini2021benchmarking}. Standard unsupervised scores---NMI, ARI, clustering accuracy after Hungarian matching~\citep{Kuhn1955Hungarian}, and neighbourhood accuracy---remain necessary for comparing subtype recovery, but they are silent on latent semantics. For shared--private models the evaluation gap is sharper still: standard leaderboards report clustering on $\mathrm{concat}(P,S)$ or on~$S$ alone, but rarely quantify whether foreign private codes predict another modality better than foreign shared codes. Without such probes, one cannot tell whether ``shared'' is shared.

The gap addressed by this thesis is therefore twofold. Methodologically, existing deep shared--private pipelines do not sufficiently constrain~$S$ and~$P$ to their intended roles when both codes are allowed to reconstruct~$X_m$ under weak independence regularisation. Evaluatively, high unsupervised subtype NMI has been over-interpreted as evidence of successful decomposition. The following chapters develop Methods~1--3---predictability constraints, joint reconstruction with VICReg, and spectral shared targets---and a probe-centred protocol to close that gap on BRCA, with KIRC reserved for a later revision.

\clearpage
\section{Background}
\label{sec:background}

This chapter provides the biological and technical background required for the methodology. Biological subsections introduce multi-omics data, TCGA cohorts, and PAM50 subtypes. Technical subsections summarise autoencoders and VAEs, contrastive learning, VICReg, spectral clustering, and the clustering and probe metrics used later.

\subsection{Multi-omics profiling in cancer}

An \emph{omics} modality is a high-dimensional molecular assay summarised as a sample-by-feature matrix. In this thesis three modalities are central: messenger RNA (mRNA) expression, DNA methylation, and microRNA (miRNA) expression. Jointly they capture transcriptional state, epigenetic regulation, and post-transcriptional control~\citep{hasin2017multi}. Because the assays are generated from the same tumour aliquots, matched multi-omics enables unsupervised integration aimed at patient subtypes rather than at single-assay artefacts~\citep{jiang2022big,mohr2024navigating}.

Challenges are substantial. Feature dimensions differ by orders of magnitude before selection; noise and missingness patterns are modality-specific; and biological signal may be shared, partially shared, or private. Linear dimension reduction such as principal component analysis remains a useful preprocessing and diagnostic tool~\citep{yeung2001principal}, but nonlinear encoders are often needed once interactions across thousands of features are considered~\citep{Goodfellow-et-al-2016}.

\subsection{TCGA BRCA and KIRC}

The Cancer Genome Atlas provides harmonised multi-omics for many tumour types. Breast invasive carcinoma (BRCA) is the primary cohort in this work. After the preprocessing used in the accompanying codebase, BRCA comprises $N=875$ samples with $1000$ mRNA features, $1000$ methylation features, and $503$ miRNA features, labelled into five PAM50 subtypes~\citep{parker2009pam50}. Kidney renal clear cell carcinoma (KIRC) is a secondary cohort with $N=289$ samples, much higher raw feature dimensionality ($58316$ mRNA, $22928$ methylation, $1879$ miRNA features in the prepared tables), and two clinical classes in the evaluation setting used here. KIRC results are deferred to a later revision; the cohort is mentioned so that the experimental design remains explicit about intended generalisation beyond BRCA.

\subsection{PAM50 breast cancer subtypes}

PAM50 is a nearest-centroid classifier based on a fifty-gene expression signature that assigns breast tumours to intrinsic subtypes such as Luminal~A, Luminal~B, HER2-enriched, Basal-like, and Normal-like~\citep{parker2009pam50}. Although derived from transcriptomics, PAM50 labels correlate with multi-omics structure and are the de facto unsupervised benchmark target for BRCA integration methods~\citep{chen2023mocss,rappoport2021evaluation}. In this thesis PAM50 labels are used only for evaluation, never for model selection: checkpoints are ranked by label-free validation loss.

\subsection{Autoencoders and variational autoencoders}

An autoencoder (AE) learns an encoder $f_{\mathrm{enc}}$ and decoder $f_{\mathrm{dec}}$ by minimising a reconstruction loss between $X$ and $\hat{X}=f_{\mathrm{dec}}(f_{\mathrm{enc}}(X))$~\citep{Goodfellow-et-al-2016}. Bottleneck latents act as compressed representations for downstream clustering. A variational autoencoder (VAE) places a probabilistic encoder $q(z\mid x)$ and a prior $p(z)$ in an evidence lower-bound objective~\citep{kingma2013auto}. OMIDIENT's private branch is a Dirichlet-oriented VAE variant rather than a standard Gaussian VAE~\citep{Safinianaini_Välimäki_Bresson_Gorbonos_Rajamäki_Aaltonen_Marttinen_2025}. Throughout Methods~1--3, unless otherwise stated, private and shared encoders follow a deterministic AE parameterisation aligned with MOCSS for fair comparison.

\subsection{Contrastive learning and InfoNCE}

Contrastive learning trains embeddings so that positive pairs are close and negatives are far in a projected space~\citep{chen2020simple}. The InfoNCE loss is a standard instantiation: for a batch of projected embeddings, each anchor is classified against negatives via a temperature-scaled softmax. In multi-omics shared--private models, positives are matched samples' shared codes across modalities. Contrastive pressure encourages cross-view alignment of~$S$, but if reconstruction also forces~$S_m$ to encode private detail, contrastive learning alone cannot purify the shared subspace.

\subsection{VICReg and independence regularisation}

Variance--Invariance--Covariance Regularization (VICReg)~\citep{bardes2022vicreg} stabilises non-contrastive learning through three terms: a variance hinge that prevents dimensional collapse, a covariance penalty that reduces redundancy within an embedding, and---when applied to two embeddings---a cross-covariance penalty that encourages linear independence. Barlow Twins pursues a related redundancy-reduction goal~\citep{zbontar2021barlow}. In this thesis a VICReg-style penalty is applied between~$P_m$ and~$S_m$ inside each modality for Method~1+VICReg and Method~2, replacing MOCSS's weak elementwise orthogonality. Linear independence is not full statistical disentanglement~\citep{martens2024disentangling}, but it is a stricter and more interpretable constraint than elementwise co-activation penalties.

\subsection{Spectral clustering and fused graphs}

Spectral clustering embeds samples using eigenvectors of a graph Laplacian and then applies $k$-means in that embedding~\citep{ng2002spectral,vonluxburg2007tutorial,lloyd1982least}. SNF-style fusion builds a single patient graph that aggregates neighbourhood agreement across modalities~\citep{wang2014snf}. Method~3 precomputes a fused $k$-nearest-neighbour graph offline, extracts a Laplacian embedding~$U\in\mathbb{R}^{N\times r}$ with $r=16$, and trains shared codes to regress toward~$U$. The NMI of $k$-means on~$U$ itself is an unsupervised ceiling for how much subtype signal the spectral target contains; on BRCA that ceiling is approximately $\mathrm{NMI}(U)\approx 0.387$ in the configuration used here.

\subsection{Clustering metrics and assignment}

Unsupervised subtype recovery is evaluated with normalised mutual information (NMI), adjusted Rand index (ARI), clustering accuracy after optimal label permutation via the Hungarian algorithm~\citep{Kuhn1955Hungarian}, and $k$-nearest-neighbour classification accuracy in the embedding space. NMI and ARI compare partitions without requiring label alignment; accuracy requires a matching between cluster indices and true labels. All clustering in this thesis uses $k$-means with $n_{\mathrm{init}}=30$~\citep{lloyd1982least}.

\subsection{Probe metrics for shared--private roles}

Because clustering metrics do not test role assignment, additional probes are required. Cross-view Ridge regression asks how well modality~$A$'s features are predicted from other modalities' shared codes versus private codes. The signed gap
\begin{equation}
\Delta\mathrm{X}\,R^{2}
\;=\;
R^{2}\!\left(X_A\leftarrow (S_B,S_C)\right)
-
R^{2}\!\left(X_A\leftarrow (P_B,P_C)\right)
\end{equation}
is positive when foreign shared codes explain~$X_A$ better than foreign private codes---the intended asymmetry. Own-view Ridge~$R^{2}$ for $X_m\leftarrow S_m$, $X_m\leftarrow P_m$, and $X_m\leftarrow[P_m;S_m]$ diagnoses collapse and redundancy within a modality. Centred kernel alignment (CKA) summarises linear similarity among shared codes across views. Embedding visualisations (UMAP) provide qualitative evidence of whether subtype geometry lives primarily in~$S$, in~$P$, or in both.

\paragraph{Interpretation rule.}
Low $\mathrm{concat}(P)$ subtype NMI is not by itself evidence of good disentanglement: private codes may have collapsed, or they may carry complementary subtype-related biology. Models are therefore judged jointly by~$\Delta\mathrm{X}\,R^{2}$, UMAP role separation, and integration NMI on $\mathrm{concat}(P,S)$ or related fused embeddings.

\clearpage
\section{Methodology}
\label{sec:methodology}

This chapter formalises shared--private multi-omics learning, describes the MOCSS and OMIDIENT baselines, introduces Methods~1--3, and defines the evaluation protocol with emphasis on~$\Delta\mathrm{X}\,R^{2}$.

\subsection{Problem formalisation}

Let $M=3$ matched omics modalities be observed for $N$ samples. For modality $m\in\{1,2,3\}$, write $X_m\in\mathbb{R}^{N\times d_m}$. A shared--private encoder produces
\begin{equation}
(S_m,P_m)=E_m(X_m),\qquad S_m,P_m\in\mathbb{R}^{N\times d},
\end{equation}
with embedding dimension $d=32$ in all experiments. Intended semantics are:
\begin{itemize}
\item $S_m$ should capture information about sample~$i$ that is predictable across modalities (cross-view biology, including subtype-related programmes);
\item $P_m$ should capture residual view-specific variation useful for reconstructing $X_m$ but not for explaining other views.
\end{itemize}
Downstream unsupervised subtyping is performed by $k$-means on $\mathrm{mean}_m(S_m)$, on $\mathrm{concat}_m(P_m)$, and on $\mathrm{concat}(P_1,P_2,P_3,\mathrm{mean}_m(S_m))$, among other probes. The scientific claim of interest is not only that some embedding clusters by subtype, but that~$S$ and~$P$ realise the decomposition above.

\paragraph{Shared implementation conventions.}
Unless a method states otherwise, encoders are four-layer multilayer perceptrons with $\tanh$ activations and widths $[512,256,128,32]$. Decoders and predictors use four-layer MLPs with $\tanh$ hidden units and a final sigmoid. Shared codes are passed through a two-layer $\tanh$ projection before contrastive learning. Weights are Kaiming-initialised with zero biases. Reconstruction losses follow the MOCSS-style centred, row-normalised matrix-norm metric rather than a literal mean squared error, ensuring fair comparison with the baselines. The contrastive term is the sum of pairwise InfoNCE losses among the three projected shared embeddings in a batch~\citep{chen2020simple,chen2023mocss}.

\subsection{Baseline: MOCSS}

MOCSS~\citep{chen2023mocss} trains, for each modality, dual reconstruction paths
\begin{equation}
\hat{X}_m^{(S)}=\mathrm{Dec}_{S}(S_m),\qquad
\hat{X}_m^{(P)}=\mathrm{Dec}_{P}(P_m),
\end{equation}
together with a weak elementwise orthogonality penalty between~$\ell_2$-normalised $S_m$ and $P_m$, and InfoNCE on $\mathrm{MLP}(S_m)$ across modalities. Dual full reconstruction is the central inductive bias under scrutiny: each code is encouraged to be individually sufficient for~$X_m$, which conflicts with a residual interpretation of~$P_m$ and invites private codes to store shared subtype signal.

\subsection{Baseline: OMIDIENT}

OMIDIENT~\citep{Safinianaini_Välimäki_Bresson_Gorbonos_Rajamäki_Aaltonen_Marttinen_2025} keeps a MOCSS-style shared autoencoder and contrastive head, but replaces the private pathway with a product-Dirichlet variational encoder--decoder and a corresponding KL regulariser. The hope is that compositional private latents better absorb view-specific residual structure. The shared branch, however, remains essentially MOCSS. This thesis therefore treats OMIDIENT as a private-branch intervention whose shared--private claims must still be verified with the same probes used for AE-private models. A Method~2 hybrid that grafts Dirichlet private codes onto Method~2's joint reconstruction objective is included as a diagnostic companion.

\subsection{Method 1: predictability surrogates}

Method~1 replaces dual solo reconstruction as the definition of shared versus private with a set of predictability constraints. For target modality~$A$ and the other two modalities $\{B,C\}$, five networks are trained:
\begin{align}
h_A(P_A,S_A)&\rightarrow X_A && \text{(within-view joint reconstruction)},\\
f_A(S_B,S_C)&\rightarrow X_A && \text{(cross-shared sufficiency)},\\
g_A(P_B,P_C)&\rightarrow X_A && \text{(cross-private adversary)},\\
k_A(S_A)&\rightarrow X_A && \text{(own-shared insufficiency)},\\
b_A(P_A)&\rightarrow X_A && \text{(own-private sufficiency)}.
\end{align}
Shared codes are additionally aligned by InfoNCE on $\mathrm{MLP}(S)$. Figure~\ref{fig:arch-m1} summarises the architecture.

\begin{figure}[htbp]
\centering
\includegraphics[width=0.98\textwidth]{M1.png}
\caption{Method~1 architecture. Per-modality encoders emit $(P_m,S_m)$. Predictors $h,f,g,k,b$ implement joint reconstruction, cross-shared sufficiency, private non-leakage, own-shared insufficiency, and private sufficiency. A contrastive head aligns projected shared codes.}
\label{fig:arch-m1}
\end{figure}

Let $\mathcal{L}_{\mathrm{std}}$, $\mathcal{L}_{\mathrm{shared}}$, $\mathcal{L}_{\mathrm{adv}}$, $\mathcal{L}_{\mathrm{own}}$, and $\mathcal{L}_{\mathrm{priv}}$ denote the sums over modalities of the corresponding reconstruction metrics. Training uses a two-step procedure per batch. Step~1 updates only the adversaries $g$ and $k$ on detached codes:
\begin{equation}
\mathcal{L}_{\mathrm{step1}}=\sum_A\bigl(\mathcal{L}_{\mathrm{adv}}^{A}+\mathcal{L}_{\mathrm{own}}^{A}\bigr).
\end{equation}
Step~2 updates encoders and constructive predictors with adversaries frozen:
\begin{equation}
\begin{aligned}
\mathcal{L}_{\mathrm{step2}}
&=
\sum_A\Bigl[
\mathcal{L}_{\mathrm{std}}^{A}
+\alpha\,\mathcal{L}_{\mathrm{shared}}^{A}
+\lambda_{\mathrm{priv}}\,\mathcal{L}_{\mathrm{priv}}^{A}
-\beta\,\mathcal{L}_{\mathrm{adv}}^{A}
-\gamma\,\mathcal{L}_{\mathrm{own}}^{A}
\Bigr]
+\lambda_{\mathrm{ctr}}\,\mathcal{L}_{\mathrm{ctr}}.
\end{aligned}
\end{equation}
The minus signs encourage encoders to make foreign private codes uninformative about~$X_A$ and to make own~$S_A$ insufficient for full reconstruction, while $b_A$ prevents the trivial collapse of~$P_A$ under privacy pressure. Checkpointing uses a constructive validation score that drops the signed adversarial terms, because raw $\mathcal{L}_{\mathrm{step2}}$ can become negative and unstable as an early-stopping signal. Across configurations, models are ranked by the label-free core validation loss $\mathcal{L}_{\mathrm{std}}+\mathcal{L}_{\mathrm{shared}}+\mathcal{L}_{\mathrm{priv}}+\mathcal{L}_{\mathrm{ctr}}$.

\subsection{Method 1 with VICReg ablation}

The incentive for this ablation is practical and diagnostic. Full Method~1 couples five predictors per modality with a two-step minimax procedure, and the own-shared adversary~$k$ (weight~$\gamma$) is particularly delicate: too strong a~$\gamma$ can strip useful shared content from~$S_A$, while too weak a~$\gamma$ leaves private bleed in~$S$. The Method~1+VICReg ablation therefore asks whether the \emph{cross-view} predictability story ($h$, $f$, $g$, and private sufficiency~$b$) already suffices when supplemented by a modern linear independence regulariser, without the brittle own-shared insufficiency term. Concretely, it removes~$k$ and adds a VICReg-style independence penalty between~$P_m$ and~$S_m$~\citep{bardes2022vicreg}. Predictors $h,f,g,b$ remain. The VICReg term combines variance hinges on both codes, within-code covariance whitening, and a cross-covariance penalty that targets linear $P\perp S$ independence. Figure~\ref{fig:arch-m1vic} illustrates the resulting architecture. Two-step training reduces to an adversary step on~$g$ only. Ranking again uses $\mathcal{L}_{\mathrm{std}}+\mathcal{L}_{\mathrm{shared}}+\mathcal{L}_{\mathrm{priv}}+\mathcal{L}_{\mathrm{ctr}}$, excluding $\mathcal{L}_{\mathrm{vic}}$ from the core rank so that independence regularisation is treated as a training regulariser rather than as the selection criterion. Empirically (Chapter~\ref{sec:results}), the ablation is more stable than full Method~1 and produces a cleaner subtype NMI ordering of~$S$ over~$P$, but it still fails the cross-view $\Delta\mathrm{X}\,R^{2}$ test---which motivates the simpler Method~2 and the spectral Method~3 that follow.

\begin{figure}[htbp]
\centering
\includegraphics[width=0.98\textwidth]{M1_VICReg.png}
\caption{Method~1+VICReg ablation. The own-shared adversary $k$ is removed; a VICReg $P\perp S$ term is added while retaining cross-view predictors $h,f,g$ and private sufficiency $b$.}
\label{fig:arch-m1vic}
\end{figure}

\subsection{Method 2: joint reconstruction and VICReg}

Method~2 abandons cross-reconstruction adversaries in favour of a simpler inductive bias. Each modality is reconstructed only from the concatenation of its own codes,
\begin{equation}
\hat{X}_m=D_m\bigl([P_m;S_m]\bigr),
\end{equation}
while VICReg encourages $P_m\perp S_m$ and InfoNCE aligns projected shared codes (Figure~\ref{fig:arch-m2}). The training objective is
\begin{equation}
\mathcal{L}
=
\lambda_{\mathrm{rec}}\,\mathcal{L}_{\mathrm{recon}}
+\lambda_{\mathrm{ctr}}\,\mathcal{L}_{\mathrm{ctr}}
+\lambda_{\mathrm{vic}}\,\mathcal{L}_{\mathrm{vic}},
\end{equation}
with all terms non-negative, so raw validation loss is safe for early stopping. Configurations are ranked by $\mathcal{L}_{\mathrm{recon}}+\mathcal{L}_{\mathrm{ctr}}$. Relative to MOCSS, Method~2 removes the requirement that~$S_m$ or~$P_m$ alone reconstruct~$X_m$, which is the most direct architectural fix of the dual-reconstruction failure mode.

\begin{figure}[htbp]
\centering
\includegraphics[width=0.98\textwidth]{M2.png}
\caption{Method~2 architecture. A single joint decoder reconstructs $X_m$ from $[P_m;S_m]$; VICReg replaces elementwise orthogonality; contrastive learning aligns shared codes.}
\label{fig:arch-m2}
\end{figure}

A companion variant, Method~2+OMIDIENT-private, replaces only the private encoder with OMIDIENT's product-Dirichlet branch while retaining Method~2's joint decoder, VICReg, and contrastive shared pathway. It is used to test whether Dirichlet private codes remain well-behaved once dual solo reconstruction is removed.

\subsection{Method 3: spectral shared targets with residual private decoding}

Method~3 supplies an explicit clustering-shaped target for shared codes. Offline, modality-wise $k$-nearest-neighbour graphs are fused and a normalised Laplacian embedding $U\in\mathbb{R}^{N\times r}$ with $r=16$ is computed~\citep{wang2014snf,ng2002spectral,vonluxburg2007tutorial}. A spectral head $f_{\mathrm{spec}}$ maps $\bar{S}=\mathrm{mean}_m(S_m)$ toward~$U$:
\begin{equation}
\mathcal{L}_{\mathrm{spec}}
=
\frac{1}{N}\sum_{i=1}^{N}
\bigl\|f_{\mathrm{spec}}(\bar{S}_i)-U_i\bigr\|_{2}^{2}.
\end{equation}
Private codes reconstruct residuals through
\begin{equation}
\hat{X}_m=\mathrm{Dec}_m\bigl(P_m,\mathrm{stopgrad}(S_m)\bigr),
\end{equation}
so gradients from residual decoding do not rewrite the shared code as a second full autoencoder of~$X_m$. Additional terms include a view-identity head on~$P$ (encouraging assay-specific content to remain in the private branch), optional VICReg, a gradient-reversal (GRL) leakage adversary that discourages private codes from predicting soft clusters of~$\bar{S}$, and a contrastive term whose weight $\lambda_{\mathrm{ctr}}$ is often selected as~$0$ under loss-based ranking. Figure~\ref{fig:arch-m3} shows the architecture. Core ranking uses $\mathcal{L}_{\mathrm{spec}}+\mathcal{L}_{\mathrm{res}}+\mathcal{L}_{\mathrm{ctr}}$. The unsupervised NMI of $k$-means on~$U$ itself, $\mathrm{NMI}(U)\approx 0.387$ on BRCA, is reported as a spectral ceiling. Full notation and formulae for the offline spectral target, spectral regression, view head, and GRL non-leakage block are collected in Appendix~\ref{app:method3}.

\begin{figure}[htbp]
\centering
\includegraphics[width=0.98\textwidth]{M3.png}
\caption{Method~3 architecture. An offline fused spectral embedding $U$ supervises mean shared codes via $f_{\mathrm{spec}}$; private decoders reconstruct residuals with stop-gradient on $S$; optional view-head and GRL leakage controls regularise role separation.}
\label{fig:arch-m3}
\end{figure}

\subsection{Training protocol}

All BRCA runs use batch size~$32$, up to~$100$ epochs, and early stopping with patience~$20$ on the method-specific validation score ($\mathrm{min\_delta}=0.005$). Learning rates and weight decays follow the MOCSS defaults used in the repository unless a sweep token indicates otherwise. Model selection is label-free: PAM50 labels never enter the ranking loss. After selection, clustering and probes are evaluated on the chosen configuration and aggregated over training seeds $\{21,22,23,24,25\}$.

\subsection{Evaluation protocol and interpretation of $\Delta\mathrm{X}\,R^{2}$}

The evaluation suite comprises:
\begin{enumerate}
\item subtype NMI, ARI, accuracy, and $k$NN accuracy on $\mathrm{mean}(S)$, $\mathrm{concat}(P)$, and $\mathrm{concat}(P,S)$;
\item cross-view $\Delta\mathrm{X}\,R^{2}$ and related Ridge probes;
\item own-view Ridge $R^{2}$ for $X\leftarrow S$, $X\leftarrow P$, and $X\leftarrow[P;S]$;
\item pairwise CKA among shared codes;
\item UMAP visualisations of $\mathrm{mean}(S)$ and $\mathrm{concat}(P)$.
\end{enumerate}

\paragraph{How to read $\Delta\mathrm{X}\,R^{2}$.}
A reliably positive gap means foreign shared codes explain another modality better than foreign private codes---evidence that cross-view signal concentrates in~$S$. Near-zero gaps indicate indistinct roles. Large positive gaps must still be checked against own-view $X\leftarrow P$: if that quantity is strongly negative, private codes have collapsed and the gap is not evidence of healthy privacy. Likewise, moderately high $\mathrm{concat}(P)$ NMI is acceptable when UMAP layouts differ from those of~$S$ and $\Delta\mathrm{X}\,R^{2}$ remains positive; private codes may then carry complementary subtype correlates rather than stolen shared geometry~\citep{levine2013integrated}.

\paragraph{Optional integration gap.}
An auxiliary quantity $\mathrm{gap}_{S-\mathrm{PS}}=\mathrm{mean}(S)\,\mathrm{NMI}-\mathrm{concat}(P,S)\,\mathrm{NMI}$ asks whether adding private codes hurts or helps the fused clustering embedding. A large positive gap means private content dilutes an already subtype-rich shared code; a negative gap means private codes contribute additional subtype structure to the concatenation. This thesis treats $\mathrm{gap}_{S-\mathrm{PS}}$ as secondary: it is useful context in compact tables, but it does not by itself certify role separation. The primary decomposition probes remain $\Delta\mathrm{X}\,R^{2}$, own-view $X\leftarrow P$, and the comparative UMAP geometries of~$S$ and~$P$.

\subsection{Strengths and limitations}

Table~\ref{tab:method-summary} summarises qualitative trade-offs. Method~1 offers the most explicit information-theoretic story but is heavy: fifteen predictors, two optimisers, and sensitive $(\alpha,\beta,\gamma,\lambda_{\mathrm{priv}})$ knobs. Method~1+VICReg reduces adversary complexity and adds linear independence, yet removing~$k$ weakens the insufficiency constraint on own~$S$. Method~2 is stable and simple, fixing MOCSS's dual-reconstruction flaw directly, but relies on VICReg and contrastive weights to prevent asymmetric collapse. Method~3 provides the strongest clustering-shaped shared target and the clearest empirical $\Delta\mathrm{X}\,R^{2}$ gains, at the cost of dependence on the quality of the offline fused graph and a spectral ceiling below the best concatenation NMI of reconstruction-based models.

\begin{table}[htbp]
\centering
\caption{Qualitative comparison of methods studied in this thesis (AE-private line unless noted).}
\label{tab:method-summary}
\begin{tabular}{p{2.4cm}p{3.2cm}p{3.2cm}p{3.6cm}}
\hline
 & Reconstruction & Independence / privacy & Main risk \\
\hline
MOCSS & dual solo $\mathrm{Dec}_S$, $\mathrm{Dec}_P$ & weak elementwise ortho & $P$ steals shared subtype \\
OMIDIENT & Dirichlet private + MOCSS shared & ortho optional; KL on $P$ & private-path collapse \\
Method~1 & $h,f,g,k,b$ predictability & adversaries + $b$ & minimax instability; weak concat NMI \\
M1+VICReg & $h,f,g,b$ + no $k$ & VICReg $P\perp S$ & low $\mathrm{concat}(P)$ $\neq$ disentangling \\
Method~2 & joint $D([P;S])$ & VICReg + contrastive & needs careful $\lambda$ balance \\
Method~3 & residual $\mathrm{Dec}(P,\mathrm{stopgrad}(S))$ & spectral $U$ + GRL leak & limited by $\mathrm{NMI}(U)$ ceiling \\
\hline
\end{tabular}
\end{table}

\clearpage
\section{Results}
\label{sec:results}

This chapter reports empirical findings on TCGA BRCA. Shared--private diagnostics are presented before subtyping leaderboards, reflecting the thesis claim that integration NMI alone is an incomplete success criterion. All headline numbers are mean~$\pm$~sample standard deviation over five training seeds $\{21,22,23,24,25\}$ unless noted. KIRC is described in the setup but numerical KIRC results are deferred to a later revision.

\subsection{Experimental setup}

\paragraph{Datasets.}
Table~\ref{tab:datasets} summarises the prepared cohorts. BRCA is the evaluation focus: three modalities after feature selection, $N=875$ samples, and five PAM50 subtypes~\citep{parker2009pam50}. KIRC is included in the repository pipeline with two evaluation classes; its results will be reported in a later revision of this thesis.

\begin{table}[htbp]
\centering
\caption{Dataset statistics for the prepared multi-omics tables used in this thesis.}
\label{tab:datasets}
\begin{tabular}{lrrrrr}
\hline
Cohort & mRNA & Methylation & miRNA & Samples & Subtypes / classes \\
\hline
BRCA & 1000 & 1000 & 503 & 875 & 5 (PAM50) \\
KIRC & 58316 & 22928 & 1879 & 289 & 2 \\
\hline
\end{tabular}
\end{table}

\paragraph{Training and selection.}
Encoders are four-layer $\tanh$ MLPs with widths $[512,256,128,32]$. BRCA models use batch size~$32$, at most~$100$ epochs, and early stopping as described in Chapter~\ref{sec:methodology}. Hyperparameter grids are method-specific; the reported ``loss-best'' row for each method is the configuration with best label-free core validation loss on the development protocol, then re-evaluated across seeds. Clustering uses $k$-means with $n_{\mathrm{init}}=30$~\citep{lloyd1982least}. Hungarian matching is used for clustering accuracy~\citep{Kuhn1955Hungarian}.

\paragraph{Compared systems.}
AE-private comparisons include MOCSS~\citep{chen2023mocss}, Method~1+$b_m$, Method~1+VICReg+$b_m$, Method~2, and Method~3. The OMIDIENT family includes OMIDIENT~\citep{Safinianaini_Välimäki_Bresson_Gorbonos_Rajamäki_Aaltonen_Marttinen_2025} and Method~2 with OMIDIENT private codes. Method~4 from the codebase is not evaluated here.

\subsection{Shared versus private analysis}

The central empirical claim is that competitive subtype NMI does not imply successful shared--private decomposition. Figures~\ref{fig:heatmap-ae}--\ref{fig:umap} and Tables~\ref{tab:shared-private-ae}--\ref{tab:shared-private-omid} support that claim.

\begin{figure}[htbp]
\centering
\includegraphics[width=\textwidth]{03_metric_heatmap_ae.png}
\caption{Compact metric heatmap for AE-private methods on BRCA (multi-seed means). Rows are methods; columns include mean$(S)$ NMI, $\mathrm{concat}(P)$ NMI, $\mathrm{concat}(P,S)$ NMI, $\Delta\mathrm{X}\,R^{2}$, own-view probes, and CKA. Colour encodes relative magnitude within each column.}
\label{fig:heatmap-ae}
\end{figure}

\begin{figure}[htbp]
\centering
\includegraphics[width=\textwidth]{03_metric_heatmap_omidient.png}
\caption{Compact metric heatmap for the OMIDIENT family on BRCA. Near-zero $\mathrm{concat}(P)$ NMI co-occurs with large $\Delta\mathrm{X}\,R^{2}$ and pathological own-view $X\leftarrow P$ scores.}
\label{fig:heatmap-omid}
\end{figure}

\begin{table}[htbp]
\centering
\caption{Shared--private diagnostic panel on BRCA (AE private). Mean~$\pm$~std over seeds $\{21,\ldots,25\}$.}
\label{tab:shared-private-ae}
\begin{tabular}{lcccc}
\hline
Method & mean$(S)$ NMI & $\mathrm{concat}(P)$ NMI & $\mathrm{concat}(P,S)$ NMI & $\Delta\mathrm{X}\,R^{2}$ \\
\hline
MOCSS & $0.418\pm0.014$ & $0.466\pm0.017$ & $0.461\pm0.008$ & $0.045\pm0.023$ \\
M1+$b_m$ & $0.414\pm0.030$ & $0.237\pm0.099$ & $0.323\pm0.063$ & $-0.038\pm0.095$ \\
M1+VICReg+$b_m$ & $0.423\pm0.021$ & $0.254\pm0.051$ & $0.394\pm0.035$ & $-0.061\pm0.071$ \\
Method~2 & $0.417\pm0.028$ & $0.389\pm0.034$ & $0.434\pm0.033$ & $0.051\pm0.044$ \\
Method~3 & $0.414\pm0.029$ & $0.443\pm0.014$ & $0.456\pm0.025$ & $0.248\pm0.023$ \\
\hline
\end{tabular}
\end{table}

\begin{table}[htbp]
\centering
\caption{Shared--private diagnostic panel on BRCA (OMIDIENT family). Mean~$\pm$~std over seeds $\{21,\ldots,25\}$.}
\label{tab:shared-private-omid}
\begin{tabular}{lcccc}
\hline
Method & mean$(S)$ NMI & $\mathrm{concat}(P)$ NMI & $\mathrm{concat}(P,S)$ NMI & $\Delta\mathrm{X}\,R^{2}$ \\
\hline
OMIDIENT & $0.426\pm0.046$ & $0.039\pm0.009$ & $0.430\pm0.054$ & $1.114\pm0.013$ \\
M2+OMID.\ private & $0.410\pm0.018$ & $0.043\pm0.002$ & $0.409\pm0.014$ & $0.770\pm0.055$ \\
\hline
\end{tabular}
\end{table}

\paragraph{MOCSS: private codes steal subtype.}
For MOCSS, $\mathrm{concat}(P)$ NMI ($0.466\pm0.017$) is essentially tied with $\mathrm{concat}(P,S)$ ($0.461\pm0.008$) and exceeds mean$(S)$ ($0.418\pm0.014$). The cross-view gap is only $\Delta\mathrm{X}\,R^{2}=0.045\pm0.023$, indicating weak preference for foreign shared codes over foreign private codes. Own-view Ridge scores reinforce the redundancy story: $X\leftarrow P$ ($0.481\pm0.006$) exceeds $X\leftarrow S$ ($0.404\pm0.013$), so private codes are highly reconstructive of the full input. Figure~\ref{fig:nmi-ae} visualises the NMI pattern; Figure~\ref{fig:umap} shows subtype geometry clearly present in MOCSS private embeddings. High unsupervised integration NMI therefore does \emph{not} imply that~$S$ and~$P$ play distinct roles.

\paragraph{OMIDIENT: collapse rather than verified privacy.}
OMIDIENT yields mean$(S)$ NMI $0.426\pm0.046$ and $\mathrm{concat}(P)$ NMI $0.039\pm0.009$, with a very large $\Delta\mathrm{X}\,R^{2}=1.114\pm0.013$. Read in isolation, these numbers look like extreme shared--private separation. Own-view probes contradict that reading: $X\leftarrow P\approx-0.342$ is strongly negative, consistent with private-path collapse of the Dirichlet branch rather than informative privacy. Single-seed shared NMI near~$0.50$ shrinks under multi-seed aggregation, cautioning against leaderboard optimism from one seed~\citep{baiao2025technical}. Method~2+OMIDIENT-private inherits the same private collapse signature ($X\leftarrow P\approx-0.341$) despite Method~2's joint decoder, showing that Dirichlet private geometry---not only MOCSS dual reconstruction---can drive the failure.

\begin{figure}[htbp]
\centering
\includegraphics[width=\textwidth]{01_nmi_bars_ae.png}
\caption{BRCA multi-seed NMI bars for AE-private methods: mean$(S)$, $\mathrm{concat}(P)$, and $\mathrm{concat}(P,S)$. MOCSS private NMI rivals concatenation; Method~1 variants suppress private NMI but also weaken concatenation; Methods~2--3 preserve competitive concatenation.}
\label{fig:nmi-ae}
\end{figure}

\begin{figure}[htbp]
\centering
\includegraphics[width=\textwidth]{01_nmi_bars_omidient.png}
\caption{BRCA multi-seed NMI bars for the OMIDIENT family. Near-zero $\mathrm{concat}(P)$ NMI should not be read as successful disentanglement without own-view probes.}
\label{fig:nmi-omid}
\end{figure}

\begin{figure}[htbp]
\centering
\includegraphics[width=\textwidth]{02_scatter_separation_vs_delta_xr2.png}
\caption{Subtype separation versus $\Delta\mathrm{X}\,R^{2}$ on BRCA. Desirable models sit toward higher separation with clearly positive $\Delta\mathrm{X}\,R^{2}$. Method~3 occupies the most favourable region among AE-private methods; OMIDIENT's extreme $\Delta$ coincides with private collapse and is not treated as a verified win.}
\label{fig:scatter}
\end{figure}

\begin{figure}[htbp]
\centering
\includegraphics[height=0.72\textheight,keepaspectratio]{05_embedding_umap.png}
\caption{UMAP of mean$(S)$ (left) and $\mathrm{concat}(P)$ (right) coloured by PAM50 subtype for the methods emphasised in this thesis (MOCSS, Method~1, Method~1+VICReg, Method~2, Method~3, OMIDIENT). Focus on those rows; any additional Method~4 row that may appear in the panel is outside the scope of this thesis and should be ignored. MOCSS private embeddings retain strong subtype structure; Method~2 shows structured but differently organised layouts on both sides; Method~3 preserves competitive shared geometry; OMIDIENT private embeddings lack subtype organisation consistent with collapse.}
\label{fig:umap}
\end{figure}

\paragraph{Method 1 and Method 1+VICReg.}
Method~1+$b_m$ reduces $\mathrm{concat}(P)$ NMI to $0.237\pm0.099$, but concatenation NMI falls to a weak $0.323\pm0.063$ and $\Delta\mathrm{X}\,R^{2}$ is noisy and on average slightly negative ($-0.038\pm0.095$). These results are not promising as a shared--private solution for two related reasons. First, the drop in $\mathrm{concat}(P)$ NMI was bought at the price of integration: the fused embedding used for subtyping becomes substantially worse than MOCSS ($0.323$ vs.\ $0.461$), so the model no longer delivers competitive unsupervised subtyping even if private subtype content has been partially suppressed. Second, the intended cross-view asymmetry still fails: $\Delta\mathrm{X}\,R^{2}$ does not become stably positive, and seed variance is large ($\pm0.095$), which is consistent with the brittleness of Method~1's two-step adversaries. In short, Method~1 can change the NMI bars without establishing that foreign~$S$ carries the shared content that foreign~$P$ does not.

Method~1+VICReg+$b_m$ is stabler and produces the most MOCSS-like \emph{look} of an AE subtype split in the sense that mean$(S)$ NMI ($0.423\pm0.021$) exceeds $\mathrm{concat}(P)$ ($0.254\pm0.051$). Yet $\Delta\mathrm{X}\,R^{2}=-0.061\pm0.071$ remains unfavourable and concatenation NMI ($0.394\pm0.035$) still trails Method~2 and Method~3. The ablation therefore improves hygiene relative to full Method~1, but still fails the probe that defines ``shared'' in the cross-view sense. The lesson is explicit: low $\mathrm{concat}(P)$ NMI is not a sufficient statistic for disentanglement.

\paragraph{Method 2 versus VICReg ablation.}
Method~2 recovers concatenation NMI $0.434\pm0.033$ with $\Delta\mathrm{X}\,R^{2}=+0.051\pm0.044$. UMAP layouts are structured on both shared and private sides but organised differently, consistent with complementary rather than duplicated geometry. This profile is preferred over Method~1+VICReg: similar shared NMI, better integration, and a positive---if modest---cross-view gap. At the same time, Method~2's $\Delta\mathrm{X}\,R^{2}$ remains close to MOCSS ($0.045$) and far below Method~3 ($0.248$). That gap is the clearest quantitative signal that Method~2 does \emph{not} yet achieve the role separation this thesis aims for. Technically, the limitation is unsurprising: VICReg only soft-penalises \emph{linear} $P$--$S$ dependence (variance hinges, within-code covariance whitening, and Frobenius cross-covariance), and the joint decoder $D([P;S])\rightarrow X$ still succeeds whenever shared and private content overlap, as long as the concatenation remains reconstructive. There is no operational pressure that forces cross-view content into~$S$ and residual content into~$P$ in the way Method~1's $f$/$g$ predictors or Method~3's spectral target do. A stronger nonlinear independence criterion, or an explicit cross-view shared objective, might close part of that gap; under the present design, Method~2 behaves like a cleaner MOCSS rather than like a verified shared--private theory. Method~3 is the method that actually delivers the desired distinction on BRCA.

\paragraph{Method 3: clearest $\Delta\mathrm{X}\,R^{2}$ win.}
Method~3 attains $\Delta\mathrm{X}\,R^{2}=0.248\pm0.023$, a clear separation from all other AE-private methods, while remaining competitive on concatenation NMI ($0.456\pm0.025$). Private concatenation NMI remains high ($0.443\pm0.014$); under the interpretation rule of Chapter~\ref{sec:methodology} this is acceptable when probes and UMAP indicate complementary biology rather than role collapse. The offline spectral target has unsupervised ceiling $\mathrm{NMI}(U)\approx0.387$, so Method~3's shared NMI ($0.414\pm0.029$) slightly exceeding that ceiling indicates that residual decoding, view heads, and encoder nonlinearities still contribute beyond pure Laplacian regression. Shared CKA for Method~3 is comparatively low ($0.547\pm0.017$ vs.\ $\approx0.76$--$0.79$ for MOCSS/Method~2/VICReg). This is expected rather than paradoxical: the loss-best Method~3 configuration uses $\lambda_{\mathrm{ctr}}=0$, so the InfoNCE alignment that usually inflates pairwise shared CKA is demoted, and the spectral objective aligns $\bar{S}$ to~$U$ without requiring the individual $S_m$ to be as linearly similar to each other as in contrastive-heavy runs. Lower CKA should therefore be read together with the strong $\Delta\mathrm{X}\,R^{2}$: Method~3 trades contrastive-driven shared homogeneity for a geometry-defined shared role.

\subsection{Subtyping performance}

Table~\ref{tab:subtyping} reports standard unsupervised subtyping metrics on the loss-best configurations. MOCSS remains strong on concatenation NMI/ARI/accuracy, as expected from a model whose private codes carry subtype signal. Method~3 matches MOCSS closely on concatenation NMI ($0.456$ vs.\ $0.461$) while delivering the superior $\Delta\mathrm{X}\,R^{2}$ already discussed. Method~2 is slightly behind on concatenation NMI but preferable to Method~1 variants that sacrifice integration for lower private NMI. OMIDIENT concatenation NMI ($0.430\pm0.054$) is competitive in mean but more variable across seeds.

\begin{table}[htbp]
\centering
\caption{Unsupervised subtyping on BRCA for loss-best configurations (mean~$\pm$~std over five seeds). Concatenation metrics use $\mathrm{concat}(P,S)$; $s_{\mathrm{NMI}}$ is NMI on mean$(S)$.}
\label{tab:subtyping}
\begin{tabular}{lcccc}
\hline
Method & concat NMI & $s_{\mathrm{NMI}}$ & ARI & Acc \\
\hline
MOCSS & $0.461\pm0.008$ & $0.418\pm0.014$ & $0.360\pm0.032$ & $0.643\pm0.060$ \\
M1+$b_m$ & $0.323\pm0.063$ & $0.414\pm0.030$ & $0.190\pm0.062$ & $0.485\pm0.060$ \\
M1+VICReg+$b_m$ & $0.394\pm0.035$ & $0.423\pm0.021$ & $0.273\pm0.070$ & $0.579\pm0.073$ \\
Method~2 & $0.434\pm0.033$ & $0.417\pm0.028$ & $0.320\pm0.049$ & $0.621\pm0.070$ \\
Method~3 & $0.456\pm0.025$ & $0.414\pm0.029$ & $0.330\pm0.047$ & $0.626\pm0.049$ \\
OMIDIENT & $0.430\pm0.054$ & $0.426\pm0.046$ & --- & --- \\
\hline
\end{tabular}
\end{table}

Own-view and alignment complements appear in Table~\ref{tab:ownview}. MOCSS shows the strongest private reconstructiveness. Method~3 balances own-view $X\leftarrow S$ and $X\leftarrow P$ more evenly while keeping a large cross-view $\Delta\mathrm{X}\,R^{2}$. OMIDIENT's negative $X\leftarrow P$ remains the decisive collapse indicator.

\begin{table}[htbp]
\centering
\caption{Own-view Ridge $R^{2}$ and shared CKA on BRCA (mean~$\pm$~std over five seeds).}
\label{tab:ownview}
\begin{tabular}{lcccc}
\hline
Method & $X\leftarrow S$ & $X\leftarrow P$ & $X\leftarrow[P;S]$ & CKA \\
\hline
MOCSS & $0.404\pm0.013$ & $0.481\pm0.006$ & $0.363\pm0.015$ & $0.761\pm0.024$ \\
M1+$b_m$ & $0.294\pm0.032$ & $0.414\pm0.028$ & $0.338\pm0.030$ & $0.589\pm0.103$ \\
M1+VICReg+$b_m$ & $0.347\pm0.040$ & $0.427\pm0.020$ & $0.346\pm0.028$ & $0.788\pm0.013$ \\
Method~2 & $0.376\pm0.013$ & $0.442\pm0.004$ & $0.377\pm0.003$ & $0.770\pm0.023$ \\
Method~3 & $0.402\pm0.006$ & $0.412\pm0.011$ & $0.375\pm0.010$ & $0.547\pm0.017$ \\
OMIDIENT & $0.427\pm0.008$ & $-0.342\pm0.000$ & $0.243\pm0.011$ & $0.757\pm0.037$ \\
\hline
\end{tabular}
\end{table}

\subsection{Method-wise findings}

\paragraph{MOCSS.}
Serves as a strong clustering baseline and as a negative example for role assignment: $\mathrm{concat}(P)\approx\mathrm{concat}(P,S)>\mathrm{mean}(S)$ with small $\Delta\mathrm{X}\,R^{2}$.

\paragraph{OMIDIENT.}
Advances private-branch modelling and remains a strong integration baseline, but near-zero $\mathrm{concat}(P)$ NMI under the probes of this thesis coincides with private-path collapse ($X\leftarrow P<0$); multi-seed evaluation and own-view probes are essential before reading that pattern as a verified shared--private win.

\paragraph{Method 1.}
Enforces predictability constraints effectively enough to shrink private NMI, but integration suffers and $\Delta\mathrm{X}\,R^{2}$ remains unstable across seeds.

\paragraph{Method 1+VICReg.}
Improves shared--private NMI ordering relative to MOCSS and stabilises CKA, yet negative $\Delta\mathrm{X}\,R^{2}$ and middling concatenation NMI make it inferior to Method~2 under the joint decision rule.

\paragraph{Method 2.}
Best simple AE-private compromise relative to Method~1 variants: joint reconstruction removes dual solo paths, VICReg replaces weak orthogonality, and multi-seed metrics remain competitive with structured UMAP on both branches. Its modest $\Delta\mathrm{X}\,R^{2}$ shows that the desired cross-view role separation is still incomplete under linear VICReg alone.

\paragraph{Method 3.}
Best probe-centred outcome: largest AE-private $\Delta\mathrm{X}\,R^{2}$, competitive concatenation NMI, and acceptable high private NMI interpreted as complementary signal. Lower shared CKA is explained by demoting contrastive alignment on the loss-best run. Performance remains bounded by spectral target quality ($\mathrm{NMI}(U)\approx0.387$).

\subsection{Hyperparameter sweeps}

Figures~\ref{fig:sweep-m1}--\ref{fig:sweep-omid} summarise label-free leaderboards over method-specific grids. Method~1 and Method~1+VICReg exhibit broad sensitivity to adversary and private-sufficiency weights: many configs that minimise private NMI also degrade concatenation NMI. Method~2 sweeps over $(\lambda_{\mathrm{rec}},\lambda_{\mathrm{ctr}},\lambda_{\mathrm{vic}})$ show a stable ridge of competitive models around moderate contrastive weight; the loss-best point used above is representative rather than an isolated spike. Method~3 sweeps indicate that larger spectral weights favour $\Delta\mathrm{X}\,R^{2}$, while $\lambda_{\mathrm{ctr}}=0$ frequently appears among loss-best settings, suggesting that the fused Laplacian target already supplies much of the cross-view alignment pressure. OMIDIENT orthogonality sweeps and Method~2+OMIDIENT-private sweeps confirm that near-zero private NMI and negative $X\leftarrow P$ persist across wide regulariser ranges, so collapse is not an artefact of a single $\lambda$.

\begin{figure}[htbp]
\centering
\includegraphics[width=\textwidth]{06_sweep_method1.png}
\caption{Hyperparameter sweep overview for Method~1 on BRCA. Performance varies sharply with predictability weights; loss-best selection does not coincide with minimal $\mathrm{concat}(P)$ NMI.}
\label{fig:sweep-m1}
\end{figure}

\begin{figure}[htbp]
\centering
\includegraphics[width=\textwidth]{06_sweep_method1_vicreg.png}
\caption{Hyperparameter sweep overview for Method~1+VICReg on BRCA.}
\label{fig:sweep-m1vic}
\end{figure}

\begin{figure}[htbp]
\centering
\includegraphics[width=\textwidth]{06_sweep_method2.png}
\caption{Hyperparameter sweep overview for Method~2 on BRCA over reconstruction, contrastive, and VICReg weights.}
\label{fig:sweep-m2}
\end{figure}

\begin{figure}[htbp]
\centering
\includegraphics[width=\textwidth]{06_sweep_method3.png}
\caption{Hyperparameter sweep overview for Method~3 on BRCA. Stronger spectral weighting improves cross-view gaps; contrastive weight is often selected as zero under core loss ranking.}
\label{fig:sweep-m3}
\end{figure}

\begin{figure}[htbp]
\centering
\includegraphics[width=\textwidth]{06_sweep_method2_omidient.png}
\caption{Sweep overview for Method~2 with OMIDIENT private codes. Private collapse signatures remain visible across the grid.}
\label{fig:sweep-m2omid}
\end{figure}

\begin{figure}[htbp]
\centering
\includegraphics[width=\textwidth]{06_sweep_omidient_ortho.png}
\caption{OMIDIENT orthogonality-weight sweep. Near-zero $\mathrm{concat}(P)$ NMI and negative own-view $X\leftarrow P$ are stable across $\lambda_{\mathrm{ortho}}$.}
\label{fig:sweep-omid}
\end{figure}

\clearpage
\section{Discussion}
\label{sec:discussion}

The experiments support a simple but often neglected claim: unsupervised subtype NMI of shared--private multi-omics models is not a certificate that~$S$ and~$P$ do what their names suggest. This chapter interprets the BRCA multi-seed evidence, relates it to prior work, and discusses limitations and implications.

\subsection{High NMI without role fidelity}

MOCSS~\citep{chen2023mocss} remains a strong clustering system on BRCA. Its failure mode is therefore scientifically informative rather than trivial. Because both shared and private decoders reconstruct the full input under only weak elementwise orthogonality, private codes are free to store subtype-discriminative biology. The empirical signature---$\mathrm{concat}(P)\approx\mathrm{concat}(P,S)>\mathrm{mean}(S)$ with $\Delta\mathrm{X}\,R^{2}\approx0.04$ and subtype-structured private UMAP---matches that architectural incentive. Related representation-learning critiques have long warned that reconstruction objectives identify compressible features, not necessarily the factors specified in a modelling narrative~\citep{bengio2013representation,martens2024disentangling}. Multi-omics shared--private models are not exempt.

The practical consequence for readers of MOCSS-style papers is interpretive caution. A table that reports concatenation NMI near~$0.46$ answers a clustering question; it does not answer whether the shared branch is the carrier of cross-omics subtype biology. In the present multi-seed audit, private codes are at least as subtype-informative as shared codes, and own-view $X\leftarrow P$ exceeds $X\leftarrow S$. That pattern is exactly what dual solo reconstruction encourages.

OMIDIENT~\citep{Safinianaini_Välimäki_Bresson_Gorbonos_Rajamäki_Aaltonen_Marttinen_2025} illustrates a complementary evaluation lesson rather than a dismissal of the method. OMIDIENT's contribution is a stronger private-branch model; its primary purpose was not to certify a new shared definition. Near-zero private subtype NMI and huge $\Delta\mathrm{X}\,R^{2}$ look like perfect separation until own-view $X\leftarrow P$ is examined. Negative reconstructiveness indicates collapse of the Dirichlet private path under the probes of this thesis, not a verified theory of shared content. Hybridising Method~2's joint decoder with OMIDIENT private codes does not remove the collapse signature, which suggests that the Dirichlet private geometry and its KL pressure---not only MOCSS dual reconstruction---can drive uninformative~$P$. The result also cautions against over-trusting single-seed leaderboards: favourable shared NMI in one run can shrink under seeds $\{21,\ldots,25\}$. Technical reproducibility studies in multi-omics make the same methodological point in broader terms~\citep{baiao2025technical,rappoport2021evaluation}.

\subsection{Why minimising private NMI is the wrong objective}

A natural but misleading response to MOCSS-style private leakage is to drive $\mathrm{concat}(P)$ subtype NMI toward zero. Method~1 and Method~1+VICReg show why that proxy fails. Both reduce private NMI relative to MOCSS, and the VICReg ablation even restores a pleasant $S>P$ NMI ordering, yet $\Delta\mathrm{X}\,R^{2}$ stays near zero or negative and concatenation NMI drops. Private codes can be made non-discriminative by collapse, by over-regularisation, or by moving information into brittle shared features that do not integrate well.

Conversely, Method~3's high private NMI is compatible with the best AE-private $\Delta\mathrm{X}\,R^{2}$ when UMAP layouts differ across branches. Assay-specific correlates of subtype are biologically real~\citep{levine2013integrated,parker2009pam50}; a private code that captures them without duplicating shared cross-view geometry is not a bug. Classical integrative factorisations already allowed individual structure to be scientifically meaningful~\citep{lock2013jive,argelaguet2018multi}. Deep models should inherit that nuance rather than treating any private subtype signal as contamination.

The decision rule used in this thesis---jointly read $\Delta\mathrm{X}\,R^{2}$, UMAP role separation, and integration NMI---is therefore not a cosmetic preference. It is required to distinguish complementary private biology from stolen shared signal and from collapsed privacy. Under that rule, Method~1+VICReg's attractive NMI split is insufficient, Method~2 is preferable among simple AE-private fixes, and Method~3 is the clearest overall win on role separation without surrendering competitive concatenation NMI.

\subsection{What the proposed methods achieve}

Method~1 operationalises shared content as cross-view predictability and private content as residual sufficiency. The idea is aligned with classical shared--individual decompositions such as JIVE and MOFA~\citep{lock2013jive,argelaguet2018multi}, translated into deep predictors $f$, $g$, $k$, and $b$. In the BRCA AE setting studied here, the adversarial estimator proved brittle: seed-to-seed variance in $\Delta\mathrm{X}\,R^{2}$ is large, and loss-best configurations under-integrate. The negative result is still useful. It shows that writing an information-theoretic story is easier than estimating it stably with fifteen networks and two-step minimax training.

Method~2 offers a more pragmatic fix: remove dual solo reconstruction, impose VICReg independence~\citep{bardes2022vicreg}, and keep contrastive alignment~\citep{chen2020simple}. Empirically it is the preferred AE-private alternative to the VICReg ablation of Method~1 when both probes and integration matter. The UMAP evidence that shared and private geometries are both structured but differently organised supports the complementary-biology reading of moderate private NMI. Method~2 also has a practical advantage for future extensions: its single-optimiser, all-positive objective is easier to hybridise than Method~1's adversaries, as the OMIDIENT-private companion experiment demonstrates even though that companion inherits private collapse. Nevertheless, Method~2's $\Delta\mathrm{X}\,R^{2}$ ($\approx0.05$) remains close to MOCSS and far below Method~3 ($\approx0.25$). That is a genuine limitation of the model, not a measurement artefact: VICReg only soft-constrains linear $P$--$S$ dependence, and joint reconstruction still permits overlapping content whenever the concatenated decoder succeeds. Without an explicit cross-view definition of sharedness, Method~2 improves hygiene relative to dual solo reconstruction but does not fully enforce the shared--private distinction sought in this thesis.

Method~3 goes further by giving shared codes an explicit spectral geometry target derived from fused patient neighbourhoods~\citep{wang2014snf,ng2002spectral,vonluxburg2007tutorial}. That choice reconnects deep shared--private learning to the classical insight that cross-view agreement lives in sample geometry, not only in per-view feature reconstruction. The large $\Delta\mathrm{X}\,R^{2}$ gain ($\approx0.25$) with competitive concatenation NMI ($\approx0.46$) is the strongest positive result of the thesis. Residual decoding with stop-gradient on~$S$ prevents the shared code from being rewritten as a second full autoencoder, while GRL leakage control and view heads provide additional role regularisation (Appendix~\ref{app:method3}). The comparatively low shared CKA of Method~3 ($\approx0.55$) is consistent with demoting contrastive alignment ($\lambda_{\mathrm{ctr}}=0$ on the loss-best run): spectral regression aligns $\bar{S}$ to~$U$ without forcing the same pairwise linear homogeneity that InfoNCE induces among $S_m$. The spectral ceiling $\mathrm{NMI}(U)\approx0.387$ clarifies a limitation and a research opportunity: Method~3 cannot conjure subtype signal absent from the fused graph, so improving graph construction, affinity kernels, or modality weighting is a direct avenue for further gains.

\subsection{Relation to evaluation standards}

Recent multi-omics benchmarking emphasises robustness, multiple metrics, and careful preprocessing disclosure~\citep{cantini2021benchmarking,rappoport2021evaluation,baiao2025technical,mohr2024navigating}. This thesis adds a missing axis for the shared--private subclass of models: role probes. Reporting only $\mathrm{mean}(S)$ NMI can hide private leakage; reporting only $\mathrm{concat}(P)$ NMI can reward collapse. Cross-view $\Delta\mathrm{X}\,R^{2}$ and own-view $X\leftarrow P$ together catch both failure modes in the BRCA results. Embedding visualisations remain indispensable qualitative checks, especially when numerical gaps are modest, as for Method~2 relative to MOCSS.

Label-free model selection deserves a brief remark. Ranking by core validation loss avoids leaking PAM50 into hyperparameter choice, which is appropriate for an unsupervised claim. It does not optimise $\Delta\mathrm{X}\,R^{2}$ directly. Some sweep cells may therefore dominate the probe metrics while losing on loss rank. The thesis treats that mismatch as an acceptable cost of honest unsupervised selection; future work could explore probe-inspired but still label-free regularisers.

\subsection{Limitations}

Several limitations bound the present conclusions. First, headline results are BRCA-only; KIRC remains to be reported in a later revision, so claims about disease transfer are provisional. Second, PAM50 labels, while standard, are transcriptomic in origin and imperfect as multi-omics ground truth~\citep{parker2009pam50,cantini2021benchmarking}. Third, VICReg enforces linear independence, not full statistical disentanglement~\citep{martens2024disentangling,zbontar2021barlow}. Fourth, Method~1's two-step adversaries may admit stronger estimators, occupancy constraints, or single-step gradient-reversal formulations not explored here. Fifth, figure panels exported from the full repository visualisation may include rows for methods outside this thesis's scope; the narrative intentionally ignores those rows. Sixth, encoder depth and width were fixed to the MOCSS-aligned MLP backbone; capacity ablations might change absolute NMI even if relative role failures persist. Finally, biological interpretation of private complementary signal remains correlative: the thesis does not map private dimensions to specific pathways or CpG modules.

\subsection{Implications for multi-omics representation learning}

For practitioners, the implication is evaluative: published shared--private models should report cross-view reconstruction gaps and own-view private reconstructiveness alongside clustering tables. A minimal reporting set suggested by this work is (i)~NMI/ARI on $\mathrm{mean}(S)$, $\mathrm{concat}(P)$, and $\mathrm{concat}(P,S)$; (ii)~$\Delta\mathrm{X}\,R^{2}$; (iii)~own-view $X\leftarrow P$; and (iv)~at least one embedding figure coloured by subtype for both branches.

For method designers, the implication is architectural. If~$P$ is meant to be residual, do not train it as a second full autoencoder of~$X$. If~$S$ is meant to be shared, give it an explicit cross-view target---predictability, contrastive agreement, or spectral geometry---that reconstruction alone does not provide. Classical factorisations already embodied these principles in linear form~\citep{lock2013jive,shen2009integrative,argelaguet2018multi}; deep models must re-enforce them rather than assume that naming latents ``shared'' and ``private'' is sufficient. Among the concrete designs studied here, Method~3 currently offers the best empirical trade-off on BRCA, with Method~2 as the simpler AE-private baseline recommended when spectral precomputation is undesirable.

\clearpage
\section{Conclusions}
\label{sec:conclusions}

This thesis examined shared--private representation learning for unsupervised multi-omics cancer subtyping, with emphasis on whether latent codes named shared ($S$) and private ($P$) actually realise that decomposition. Using TCGA BRCA multi-omics data and multi-seed evaluation over seeds $\{21,\ldots,25\}$, the work audited MOCSS and OMIDIENT, proposed Methods~1--3 that enforce clearer role separation, and ranked models by a probe-centred protocol rather than by private-subtype NMI alone.

The first conclusion is negative but central: high unsupervised subtype NMI does not imply faithful shared--private roles. In MOCSS, private concatenation NMI approximately matches full concatenation ($\approx0.47$ vs.\ $\approx0.46$) and exceeds mean shared NMI ($\approx0.42$), while $\Delta\mathrm{X}\,R^{2}$ remains near~$0.04$ and UMAP shows subtype structure in private embeddings. In OMIDIENT, near-zero private NMI and large $\Delta\mathrm{X}\,R^{2}$ ($\approx1.11$) coexist with collapsed own-view private reconstructiveness ($X\leftarrow P\approx-0.34$), so the apparent separation is not evidence of a verified shared theory. Single-seed optimism further shrinks under multi-seed aggregation.

The second conclusion is methodological. Models should be judged jointly by cross-view $\Delta\mathrm{X}\,R^{2}$, qualitative UMAP role separation, and integration NMI. Minimising $\mathrm{concat}(P)$ subtype NMI is not a valid primary objective: Method~1+$b_m$ and Method~1+VICReg achieve lower private NMI (approximately~$0.24$--$0.25$) without securing stable positive $\Delta\mathrm{X}\,R^{2}$ or competitive concatenation performance, whereas Method~3 tolerates high private NMI ($\approx0.44$) while delivering the clearest AE-private cross-view gap.

The third conclusion is constructive. Among the proposed approaches, Method~3 provides the strongest empirical support for improved shared--private distinction on BRCA ($\Delta\mathrm{X}\,R^{2}\approx0.25$ with concatenation NMI$\approx0.46$), subject to the spectral ceiling of the offline fused Laplacian target ($\mathrm{NMI}(U)\approx0.387$). Method~2 is the preferred simpler AE-private alternative to Method~1+VICReg, combining joint reconstruction, VICReg independence, and competitive multi-seed integration (concatenation NMI$\approx0.43$, $\Delta\mathrm{X}\,R^{2}\approx+0.05$) with structured but non-identical shared and private geometries.

KIRC experiments are reserved for a later revision. Future work should improve fused-graph targets for Method~3, harden predictability estimators beyond Method~1's two-step adversaries, test whether non-collapsing private priors can be combined with Method~2-style joint reconstruction, and extend the probe suite to additional tumour cohorts and omics combinations. More broadly, deep multi-omics models that advertise shared--private structure should treat that structure as an empirical claim to be tested, not as a consequence of architecture diagrams or clustering leaderboards.


\clearpage

\thesisbibliography
\bibliographystyle{apalike}
\bibliography{bibliography.bib}

\clearpage
\thesisappendix
\clearpage
\section{Method~3 spectral details}
\label{app:method3}

This appendix expands the Method~3 construction summarised in Section~\ref{sec:methodology}. Notation matches the main text: modalities $m\in\{1,2,3\}$, encoders produce $(S_m,P_m)\in\mathbb{R}^{d}$ with $d=32$, and $\bar{S}=(S_1+S_2+S_3)/3$.

\subsection{Offline fused spectral target}

For each modality~$m$, let $X_m\in\mathbb{R}^{N\times d_m}$ denote the feature matrix of the $N$ patients used to build the graph (unlabelled; no PAM50 information). Apply PCA fitted on the training split to $d_{\mathrm{pca}}=50$ dimensions, then $\ell_2$-normalise each patient row. Build a mutual $k$-nearest-neighbour affinity with $k=15$ and cosine similarity:
\begin{equation}
A_m[i,j]
=
\begin{cases}
\cos(z_{m,i},z_{m,j}), & j\in\mathrm{kNN}_k(i)\ \text{or}\ i\in\mathrm{kNN}_k(j),\\
0, & \text{otherwise},
\end{cases}
\end{equation}
symmetrise by $A_m\leftarrow\max(A_m,A_m^\top)$, and zero the diagonal. Fuse the three affinities by averaging,
\begin{equation}
A_{\mathrm{fused}}=\frac{1}{3}\bigl(A_1+A_2+A_3\bigr).
\end{equation}
Let $D=\mathrm{diag}\bigl(\sum_j A_{\mathrm{fused}}[i,j]\bigr)$ and form the normalised Laplacian
\begin{equation}
L = I - D^{-1/2} A_{\mathrm{fused}} D^{-1/2}.
\end{equation}
Compute the eigendecomposition $L\,U_{\mathrm{full}}=U_{\mathrm{full}}\Lambda$ and retain the $r=16$ eigenvectors corresponding to the smallest eigenvalues \emph{after} the trivial (near-constant) eigenvector. Row-normalise the resulting matrix $U\in\mathbb{R}^{N\times r}$. The matrix~$U$ is the offline spectral target used throughout Method~3 training. The unsupervised NMI of $k$-means on~$U$ itself is reported as the spectral ceiling $\mathrm{NMI}(U)$.

\subsection{Spectral regression}

A two-layer head $f_{\mathrm{spec}}:\mathbb{R}^{d}\rightarrow\mathbb{R}^{r}$ (hidden $\tanh$, linear output) maps the mean shared code to the spectral target:
\begin{equation}
\hat{U}_i = f_{\mathrm{spec}}(\bar{S}_i),
\qquad
\mathcal{L}_{\mathrm{spec}}
=
\frac{1}{N}\sum_{i=1}^{N}
\bigl\|\hat{U}_i - U_i\bigr\|_{2}^{2}.
\end{equation}
This term is the operational definition of sharedness in Method~3:~$S$ is trained to match multi-omics neighbourhood geometry rather than to reconstruct assay intensities.

\subsection{Residual private reconstruction}

For each modality,
\begin{equation}
\hat{X}_m
=
\mathrm{Dec}_m\bigl(\mathrm{concat}(P_m,\mathrm{stopgrad}(S_m))\bigr),
\qquad
\mathcal{L}_{\mathrm{res}}
=
\sum_{m=1}^{3}
\mathrm{recon}(\hat{X}_m,X_m),
\end{equation}
where $\mathrm{recon}$ is the MOCSS-style centred, row-normalised matrix-norm metric. Stop-gradient on~$S_m$ ensures that residual decoding cannot redefine the shared code as a second full autoencoder of~$X_m$.

\subsection{View-identity head}

A classifier $h_{\mathrm{view}}:\mathbb{R}^{d}\rightarrow\mathbb{R}^{3}$ is applied to each private code:
\begin{equation}
\mathcal{L}_{\mathrm{view}}
=
\sum_{m=1}^{3}
\mathrm{CE}\bigl(h_{\mathrm{view}}(P_m),\, m-1\bigr).
\end{equation}
Minimising $\mathcal{L}_{\mathrm{view}}$ encourages~$P_m$ to retain assay-identifying information, i.e.\ content that is useful for recognising the modality and therefore less likely to be purely shared biology.

\subsection{GRL non-leakage adversary}

Let $q_i$ be a soft cluster assignment derived from stop-gradient features of~$\bar{S}$ (refreshed cluster prototypes / soft labels with $K$ clusters). An adversary $h_{\mathrm{adv}}$ is applied to private codes through a gradient-reversal layer $\mathrm{GRL}_\lambda$:
\begin{equation}
\mathcal{L}_{\mathrm{leak}}
=
\sum_{m=1}^{3}
\mathrm{CE}\bigl(h_{\mathrm{adv}}(\mathrm{GRL}_\lambda(P_m)),\, q\bigr).
\end{equation}
Forward evaluation trains $h_{\mathrm{adv}}$ to predict~$q$ from~$P_m$; reversed gradients encourage the encoder to make~$P_m$ uninformative about the shared soft clustering. This is the non-leakage regulariser that complements residual reconstruction.

\subsection{Full Method~3 objective}

Optional VICReg and demoted contrastive terms may be included. The training objective used in the BRCA sweeps is
\begin{equation}
\begin{aligned}
\mathcal{L}
&=
\lambda_{\mathrm{spec}}\,\mathcal{L}_{\mathrm{spec}}
+\lambda_{\mathrm{res}}\,\mathcal{L}_{\mathrm{res}}
+\lambda_{\mathrm{ctr}}\,\mathcal{L}_{\mathrm{ctr}}
+\lambda_{\mathrm{vic}}\,\mathcal{L}_{\mathrm{vic}}
+\lambda_{\mathrm{view}}\,\mathcal{L}_{\mathrm{view}}
+\lambda_{\mathrm{leak}}\,\mathcal{L}_{\mathrm{leak}}.
\end{aligned}
\end{equation}
Label-free ranking uses the core score $\mathcal{L}_{\mathrm{spec}}+\mathcal{L}_{\mathrm{res}}+\mathcal{L}_{\mathrm{ctr}}$ (unweighted), treating VICReg, view, and leak terms as training regularisers. The loss-best BRCA configuration reported in the main text uses $\lambda_{\mathrm{ctr}}=0$, which helps explain the comparatively low shared CKA of Method~3.


\end{document}
